\subsection{ADJOINT REPRESENTATION}

Let \(G\) be a Lie group. Consider the law of analytic left operation

\[
(g, g ^ {\prime}) \mapsto g g ^ {\prime} g ^ {- 1} = (\operatorname{Int} g) g ^ {\prime}
\]

of G into G. This law of operation defines, by no. 3, a bilinear mapping of \(\mathcal{T}^{(\infty)}(\mathrm{G})\times \mathcal{T}^{(\infty)}(\mathrm{G})\) into \(\mathcal{T}^{(\infty)}(\mathrm{G})\) , which we shall denote by \(\top\) in this no. By Proposition 13 of no. 3,

\[
(t * t ^ {\prime}) \top t ^ {\prime \prime} = t \top (t ^ {\prime} \top t ^ {\prime \prime})\tag{22}
\]

for all \(t, t', t''\) in \(\mathcal{T}^{(\infty)}(\mathrm{G})\) . By Proposition 14 (i) of no. 3,

\[
\varepsilon_ {g} \top t = (\operatorname{Int} g) _ {*} t\tag{23}
\]

for all \(g \in G\) and \(t \in \mathcal{T}^{(\infty)}(G)\) . In particular, the mapping \(t \mapsto \varepsilon_{g} \top t\) of \(\mathcal{T}^{(\infty)}(G)\) into \(\mathcal{T}^{(\infty)}(G)\) is an automorphism of the bigebra \(\mathcal{T}^{(\infty)}(G)\) . Its restrictions to \(U(G)\) , \(U_{s}(G)\) , \(L(G)\) are denoted by \(\mathrm{Ad}_{\mathrm{U(G)}}(g)\) , \(\mathrm{Ad}_{\mathrm{U_s(G)}}(g)\) , \(\mathrm{Ad}_{\mathrm{L(G)}}(g)\) . We often write \(\mathrm{Ad}(g)\) instead of \(\mathrm{Ad}_{\mathrm{L(G)}}(g)\) when no confusion is possible. By (23), \(\mathrm{Ad}(g)\) is the tangent mapping at \(e\) to \(\operatorname{Int}(g)\) . It is an automorphism of the normable

Lie algebra \(\mathbf{L}(\mathbf{G})\) . When \(\mathbf{K}\) is of characteristic 0, \(\operatorname{Ad}_{\mathbf{U}(\mathbf{G})}(g)\) is the unique automorphism of \(\mathbf{U}(\mathbf{G})\) which extends \(\operatorname{Ad}(g)\) .

If \(\phi\) is a morphism of the Lie group G into a Lie group H, then

\[
\phi_ {*} (t \top t ^ {\prime}) = \phi_ {*} (t) \top \phi_ {*} (t ^ {\prime})\tag{24}
\]

for all \(t, t'\) in \(\mathcal{T}^{(\infty)}(\mathrm{G})\) ; this follows from Proposition 15 of no. 3.

PROPOSITION 43. Let \(t, u\) be in \(\mathcal{T}^{(\infty)}(\mathrm{G})\) . Let \(\sum_{i=1}^{n} t_i \otimes t'_i\) be the image of \(t\) under the coproduct. Then

\[
t \top u = \sum_ {i = 1} ^ {n} t _ {i} * u * t _ {i} ^ {\prime \vee}.
\]

By definition, \(t \top u\) is the image of \(t \otimes u\) under the mapping \((g, g') \mapsto gg'g^{-1}\) of \(G \times G\) into \(G\) . Now this mapping is obtained by composing the following mappings:

\[
\begin{array}{l l} \alpha \colon (g, g ^ {\prime}) \mapsto (g, g, g ^ {\prime}) & \text {of G \times G into G \times G \times G} \\ \beta \colon (g, g ^ {\prime}, g ^ {\prime \prime}) \mapsto (g, g ^ {\prime - 1}, g ^ {\prime \prime}) & \text {of G \times G \times G into G \times G \times G} \\ \gamma \colon (g, g ^ {\prime}, g ^ {\prime \prime}) \mapsto g g ^ {\prime \prime} g ^ {\prime} & \text {of G \times G \times G into G}. \end{array}
\]

On the other hand:

\[
\begin{array}{c} \alpha_ {*} (t \otimes u) = \sum_ {i = 1} ^ {n} (t _ {i} \otimes t _ {i} ^ {\prime}) \otimes u = \sum_ {i = 1} ^ {n} t _ {i} \otimes t _ {i} ^ {\prime} \otimes u \\ \beta_ {*} (\sum_ {i = 1} ^ {n} t _ {i} \otimes t _ {i} ^ {\prime} \otimes u) = \sum_ {i = 1} ^ {n} t _ {i} \otimes t _ {i} ^ {\prime \vee} \otimes u \\ \gamma_ {*} (\sum_ {i = 1} ^ {n} t _ {i} \otimes t _ {i} ^ {\prime \vee} \otimes u) = \sum_ {i = 1} ^ {n} t _ {i} * u * t _ {i} ^ {\prime \vee}. \end{array}
\]

COROLLARY 1. Let \(u \in \mathrm{L}(\mathbf{G})\) and \(u' \in \mathcal{T}^{(\infty)}(\mathbf{G})\) . Then \(u \top u' = u * u' - u' * u\) .

The image of u under the coproduct is \(u \otimes \varepsilon_{e} + \varepsilon_{e} \otimes u\) , whence

\[
u \top u ^ {\prime} = u * u ^ {\prime} * \varepsilon_ {e} + \varepsilon_ {e} * u ^ {\prime} * u ^ {\vee} = u * u ^ {\prime} - u ^ {\prime} * u.
\]

COROLLARY 2. Let \(t \in \mathcal{T}^{(\infty)}(\mathrm{G})\) and \(g \in \mathrm{G}\) . Then \(\varepsilon_g \top t = \varepsilon_g * t * \varepsilon_g^{-1}\) . If \(t \in \mathrm{L}(\mathrm{G})\) , then \(\varepsilon_g \top t = gtg^{-1}\) (where the latter product is evaluated in the group \(\mathrm{T}(\mathrm{G})\) ).

The image of \(\varepsilon_{g}\) under the coproduct is \(\varepsilon_{g} \otimes \varepsilon_{g}\) .

COROLLARY 3. Let \(a \in \mathrm{L}(\mathrm{G})\) . The vector field defined by \(a\) and the left operation \(g \mapsto \operatorname{Int} g\) of \(\mathrm{G}\) on \(\mathrm{G}\) is the field \(\mathbf{R}_a - \mathbf{L}_a\) .

The value of this field at \(g\) is

\[
\begin{array}{r l} a \top \varepsilon_ {g} = a * \varepsilon_ {g} - \varepsilon_ {g} * a & (\text { Corollary   1 }) \\ = (R _ {a}) _ {g} - (L _ {a}) _ {g} & (\text { Definition   5 }). \end{array}
\]

For all \(g \in G\) and all \(t \in \mathrm{L}(\mathrm{G})\) ,

\[
(\mathrm{Ad} g) (t) = \varepsilon_ {g} \top t = \varepsilon_ {g} * t * \varepsilon_ {g ^ {- 1}} = g t g ^ {- 1}.\tag{25}
\]

Since Ad \(g = \mathrm{T}_{e}(\mathrm{Int} g)\) , Proposition 42 of no. 11 proves that Ad is an analytic linear representation of G on the normable space \(\mathrm{L}(\mathrm{G})\) .

DEFINITION 8. The representation Ad of G on L(G) is called the adjoint representation of G.

PROPOSITION 44. For all \(a \in \mathbf{L}(\mathbf{G})\) ,

\[
(\mathrm{L} (\mathrm{Ad})) (a) = \mathrm{ad} _ {\mathrm{L} (\mathrm{G})} a.
\]

Let \(b \in \mathrm{L}(\mathrm{G})\) . By Proposition 42 (ii) of no. 11 and Corollary 3 to Proposition 43,

\[
(\mathrm{L} (\mathrm{Ad})) (a). b = - [ \mathrm{R} _ {a} - \mathrm{L} _ {a}, \mathrm{L} _ {b} ] (e).
\]

Now \(R_{a} \circ L_{b} = L_{b} \circ R_{a}\) (no. 6, Proposition 23 (ii)), whence \([R_{a}, L_{b}] = 0\) ; then, using Proposition 23 (ii),

\[
(\mathrm{L} (\mathrm{Ad})) (a). b = [ \mathrm{L} _ {a}, \mathrm{L} _ {b} ] (e) = \mathrm{L} _ {[ a, b ]} (e) = [ a, b ] = (\mathrm{ad} _ {\mathrm{L} (\mathrm{G})} a) b.
\]

PROPOSITION 45. Suppose that G is finite-dimensional and that K is of characteristic 0. Let s be an integer \(\geqslant0\) . Then the mapping \(\pi: g \mapsto \operatorname{Ad}_{\mathrm{U}_{s}(\mathbf{G})}(g)\) is an analytic linear representation of G on \(\mathrm{U}_{s}(\mathrm{G})\) and \(\mathrm{L}(\pi)a = \operatorname{ad}_{\mathrm{U}_{s}(\mathbf{G})}a\) for all \(a \in \mathrm{L}(\mathrm{G})\) .

The linear representation \(\pi\) is a quotient of \(\bigoplus_{r=0}^{\infty} \mathrm{T}^{r}(\mathrm{Ad})\) and is hence analytic. For \(a \in \mathrm{L}(\mathrm{G})\) and \(x_{1}, x_{2}, \ldots, x_{s}\) in \(\mathrm{L}(\mathrm{G})\) ,

\[
\begin{array}{l l} (\mathrm{L} (\pi) a) (x _ {1} x _ {2} \dots x _ {s}) = \sum_ {i = 1} ^ {s} x _ {1} \dots (\mathrm{L} (\mathrm{Ad}) a. x _ {i}) \dots x _ {s} & \text {(Proposition 41)} \\ = \sum_ {i = 1} ^ {s} x _ {1} \dots ([ a, x _ {i} ]) \dots x _ {s} & \text {(Proposition 44)} \\ = (\mathrm{Ad} _ {\mathrm{U} _ {s} (\mathrm{G})} a) (x _ {1} x _ {2} \dots x _ {s}). \end{array}
\]

PROPOSITION 46. Let \(h \in \mathbf{G}\) , \(x \in \mathrm{T}_h(\mathbf{G})\) and \(a \in \mathrm{L}(\mathbf{G})\) . Let \(\phi\) be the mapping \((g, g') \mapsto gg'g^{-1}\) of \(\mathbf{G} \times \mathbf{G}\) into \(\mathbf{G}\) . The image \(y\) of \((a, x) \in \mathrm{T}_e(\mathbf{G}) \times \mathrm{T}_h(\mathbf{G})\) under \(\mathrm{T}_{(e,h)}(\phi)\) is \(y = x + h((\mathrm{Ad}h^{-1})a - a)\) .

We have

\[
\begin{array}{r l} & y = (\mathrm{T} _ {(e, h)} \phi) (a \otimes \varepsilon_ {h} + \varepsilon_ {e} \otimes x) \\ & \quad = a \top \varepsilon_ {h} + \varepsilon_ {e} \top x \\ & \quad = a * \varepsilon_ {h} - \varepsilon_ {h} * a + x \\ & \quad = h ((\mathrm{Ad} h ^ {- 1}) a) - h a + x. \end{array}
\]

PROPOSITION 47. Let G be a Lie group, H and E Lie subgroups of G and suppose that \(hEh^{-1} = E\) for all \(h \in H\) . Then \(\mathcal{T}^{(\infty)}(H) \top \mathcal{T}^{(\infty)}(E) \subset \mathcal{T}^{(\infty)}(E)\) . In particular, \(\operatorname{Ad}(H)(L(E)) \subset L(E)\) and \([L(H), L(E)] \subset L(E)\) .

If \(t \in \mathcal{T}^{(\infty)}(\mathrm{H})\) and \(t' \in \mathcal{T}^{(\infty)}(\mathrm{E})\) , then \(t \otimes t' \in \mathcal{T}^{(\infty)}(\mathrm{H} \times \mathrm{E})\) and the image of \(\mathrm{H} \times \mathrm{E}\) under the mapping \((g, g') \mapsto gg'g^{-1}\) is contained in \(\mathrm{E}\) .

PROPOSITION 48. Let G be a Lie group and H and E Lie subgroups of G. Suppose that G is, as a Lie group, the semi-direct product of H by E. Let \(\rho\) be the linear representation \(g \mapsto (\mathrm{Ad} g) \mid \mathrm{L}(E)\) of the Lie group G on \(\mathrm{L}(E)\) (cf.~Proposition 47) and let \(\sigma\) be the restriction of \(\rho\) to H. Then:

\begin{enumerate}
\def\labelenumi{(\roman{enumi})}
\item
  \(\mathbf{L}(\mathbf{G})\) is the topological direct sum of \(\mathbf{L}(\mathbf{H})\) and \(\mathbf{L}(\mathbf{E})\) ;
\item
  L(H) is a subalgebra of L(G) and L(E) is an ideal of L(G);
\item
  \(\mathbf{L}(\sigma)\) is a linear representation of \(\mathbf{L}(\mathbf{H})\) on the Lie algebra of derivations of \(\mathbf{L}(\mathbf{E})\) ;
\item
  L(G) is the semi-direct product of L(H) by L(E) defined by L(σ) (Chapter I, § 1, no. 8).
\item
  is obvious and (ii) follows from Proposition 47. \(\mathbf{L}(\sigma) = \mathbf{L}(\rho) \mid \mathbf{L}(\mathbf{H})\) . Now by Propositions 40 (no. 11) and 44 (no. 12), \(\mathbf{L}(\rho)(t)\) is, for all \(t \in \mathbf{L}(\mathbf{G})\) , the restriction of \(\mathrm{ad}_{L(\mathbf{G})} t\) to \(\mathbf{L}(\mathbf{E})\) . This proves (iii). Using (i) and (ii), this also proves (iv).
\end{enumerate}

COROLLARY. Let G be a Lie group. Let \(\mathrm{T}_{e}(\mathrm{G})\) be given its unique commutative Lie algebra structure. Let \(\tau\) be an adjoint representation of \(\mathrm{L}(\mathrm{G})\) . Then the Lie algebra of \(\mathrm{T}(\mathrm{G})\) is the semi-direct product of \(\mathrm{L}(\mathrm{G})\) by \(\mathrm{T}_{e}(\mathrm{G})\) defined by \(\tau\) . In other words, for \(x, x'\) in \(\mathrm{L}(\mathrm{G})\) and \(y, y'\) in \(\mathrm{T}_{e}(\mathrm{G})\) ,

\[
[ (x, y), (x ^ {\prime}, y ^ {\prime}) ] = ([ x, x ^ {\prime} ], [ x, y ^ {\prime} ] + [ y, x ^ {\prime} ])
\]

(where the bracket on the left is evaluated in \(\mathbf{L}(\mathbf{T}(\mathbf{G}))\) and the brackets on the right in \(\mathbf{L}(\mathbf{G})\) ).

This follows from Proposition 48 and Proposition 6 of § 2, no. 2.

PROPOSITION 49. Let A be a complete normable unital associative algebra. We identify A with L(A*). Then, if \(g \in \mathbf{A}^*\) and \(y \in \mathbf{A}\) , (Ad g) \(y = gyg^{-1}\) .

Recall that \(\operatorname{Ad} g = \mathrm{T}_{1}(\operatorname{Int} g)\) . Let \(u_{g}\) be the mapping \(x \mapsto gxg^{-1}\) of A into A. The identity chart of \(\mathrm{A}^{*}\) into A transforms \(\operatorname{Int} g\) into \(u_{g} \mid \mathrm{A}^{*}\) . The tangent mapping at each point of \(\mathrm{A}^{*}\) to this mapping is equal to \(u_{g}\) , whence the proposition.

COROLLARY. For all \(g \in \mathbf{A}^*\) , let \(i(g)\) be the automorphism \(y \mapsto \mathrm{gyg}^{-1}\) of \(\mathbf{A}\) , so that \(i\) is an analytic linear representation of \(\mathbf{A}^*\) on \(\mathbf{A}\) . For all \(z \in \mathrm{L}(\mathbf{A}^*) = \mathbf{A}\) , \(\mathrm{L}(i)z\) is the inner derivation \(y \mapsto zy - yz\) of \(\mathbf{A}\) .

This follows from Propositions 49 and 44.

\subsection{TENSORS AND INVARIANT FORMS}

Let G be a Lie group. We consider G as operating on itself by left (resp. right) translation. Let \(\lambda\) be a vector functor of class \(C^{\omega}\) for isomorphisms. Then \(\lambda(\mathrm{TG})\) is an analytic left (resp. right) vector G-bundle (§1, no. 8, Corollary to Proposition 16). The mapping \((g, u) \mapsto gu\) (resp. ug) of \(G \times \lambda(\mathrm{L}(G))\) onto \(\lambda(\mathrm{TG})\) is an isomorphism \(\phi\) (resp. \(\psi\) ) of vector G-bundles (§7, no. 8, Corollary 2 to Proposition 17). Every G-invariant section of \(\lambda(\mathrm{TG})\) is analytic and determined by its value at e (§1, no. 8, Corollary 1 to Proposition 17). Such a section is called left (resp. right) invariant. Let \(\sigma\) be a left invariant section of \(\lambda(\mathrm{TG})\) ; the transform \(\sigma'\) of \(\sigma\) under a right translation \(\delta(g)\) is defined by \(\sigma'(\delta(g)h) = \lambda(\mathrm{T}_{h}(\delta(g)))\sigma(h)\) for all \(h \in G\) ; it is also left invariant; it is also derived from \(\sigma\) by \(\gamma(g) \circ \delta(g) = \operatorname{Int}(g)\) and hence

\[
\sigma^ {\prime} (e) = \lambda (\mathrm{Ad} g). \sigma (e).\tag{26}
\]

Similarly, let \(\tau\) be a right invariant section of \(\lambda(\mathrm{TG})\) ; the transform \(\tau'\) of \(\tau\) under a left translation \(\gamma(g)\) is also right variant and

\[
\tau^ {\prime} (e) = \lambda (\mathrm{Ad} g). \tau (e).\tag{27}
\]

We now consider G × G as operating on G on the left by

\[
\left(\left(g, g ^ {\prime}\right), g ^ {\prime \prime}\right) \mapsto g g ^ {\prime \prime} g ^ {\prime - 1}.
\]

Then G is a left Lie homogeneous space of \(G \times G\) (§ 1, no. 6, Example). Hence \(\lambda(\mathrm{TG})\) is an analytic left vector \((\mathrm{G} \times \mathrm{G})\) -bundle. A section of \(\lambda(\mathrm{TG})\) is called biinvariant if it is invariant under the action of \(\mathrm{G} \times \mathrm{G}\) on \(\lambda(\mathrm{TG})\) , in other words if it is invariant under left and right translations. Let \(\lambda(\mathrm{L}(\mathrm{G}))_{0}\) be the set of elements of \(\lambda(\mathrm{L}(\mathrm{G}))\) invariant under \(\lambda(\mathrm{Ad}(\mathrm{G}))\) . For all \(u \in \lambda(\mathrm{L}(\mathrm{G}))_{0}\) , let \(\sigma_{u}\) be the mapping of G into \(\lambda(\mathrm{TG})\) defined by \(\sigma_{u}(g) = gu = ug\) . Then \(u \mapsto \sigma_{u}\) is a bijection of \(\lambda(\mathrm{L}(\mathrm{G}))_{0}\) onto the set of biinvariant sections of \(\lambda(\mathrm{TG})\) (§ 1, no. 8, Corollary 1 to Proposition 17).

PROPOSITION 50. Let G be a Lie Group (assumed to be finite-dimensional if K is of characteristic \textgreater0). Let E be the vector space of continuous alternating multilinear forms of degree k on \(T_{e}(G)\) . For all \(u \in E\) , let \(\omega^{u}\) be the differential form of degree k on G such that \((\omega^{u})_{g}\) is the multilinear form on \(T_{e}(G)\) derived from u by the translation \(h \mapsto gh\) (resp. \(h \mapsto hg\) ). Then \(\omega^{u}\) is analytic and left (resp. right) invariant on G. The mapping \(u \mapsto \omega^{u}\) is an isomorphism of E onto the vector space of left (resp. right) invariant differential forms of degree k on G.

This is a special case of what we have said above.

Let \(\mathbf{F}\) be a complete normable space. Proposition 50 remains true if differential forms on \(\mathbf{G}\) with values in \(\mathbf{K}\) are replaced by differential forms on \(\mathbf{G}\) with values in \(\mathbf{F}\) . For every continuous linear mapping \(u\) of \(\mathrm{T}_e(\mathrm{G})\) into \(\mathbf{F}\) , there exists a differential form \(\omega^u\) of degree 1 on \(\mathbf{G}\) , with values in \(\mathbf{F}\) , such that \((\omega^{u})_{g} = u \circ \mathrm{T}_{g}(\gamma(g)^{-1})\) . In particular, take \(\mathrm{F} = \mathrm{T}_e(\mathrm{G})\) and \(u = \mathrm{Id}_{\mathrm{T_e(G)}}\) . We then obtain the differential form \(\omega\) on G such that \(\omega_{g} = \mathrm{T}_{g}(\gamma (g^{-1}))\) ; this differential form is left invariant and analytic; it is called the left canonical differential form of G. \(\omega_{g}(t) = g^{-1}t\) for all \(t\in \mathrm{T}_g(\mathrm{G})\) .

If F is again an arbitrary complete normable space and \(u \in \mathcal{L}(\mathrm{T}_{e}(\mathrm{G}), \mathrm{F})\) , then \(\omega^{u} = u \circ \omega\) . In particular (taking F = K), the mapping \(v \mapsto v \circ \omega\) is a linear bijection of the dual of \(\mathrm{T}_{e}(\mathrm{G})\) onto the vector space of differential forms of degree 1 with values in K which are left invariant under G.

Similarly, the differential form \(\omega'\) on G such that \(\omega_g' = \mathrm{T}_g(\delta(g))\) is called the right canonical differential form of G. There are analogous properties to those of \(\omega\) , which we leave to the reader to state. The mapping \(g \mapsto g^{-1}\) of G onto G transforms \(\omega\) into \(\omega'\) .

\subsection{MAURER-CARTAN FORMULAE}

Let X be a manifold of class \(C^{r}\) , of finite dimension if K is of characteristic \textgreater0, and let L be a complete normable Lie algebra. Let \(\alpha\) be a differential form of degree 1 on X with values in L of class \(C^{r-1}\) . Let \(x \in X\) . The mapping

\[
(u _ {1}, u _ {2}) \mapsto [ \alpha_ {x} (u _ {1}), \alpha_ {x} (u _ {2}) ]
\]

of \(\mathrm{T}_x(\mathbf{X})\times \mathrm{T}_x(\mathbf{X})\) into L is a continuous alternating bilinear form on \(\mathrm{T}_x(\mathbf{X})\) with values in L. We shall denote it by \([\alpha ]_x^2\) so that \([\alpha ]^2\) is a differential form of degree 2 on X with values in L. Identifying an open neighbourhood of \(x\) in X with an open subset of a Banach space, we see immediately that \([\alpha ]^2\) is of class \(\mathrm{C}^{r - 1}\) . If \(\mathbf{X}'\) is a manifold of class \(\mathrm{C}^r\) and \(f\colon \mathbf{X}'\to \mathbf{X}\) is a morphism, then

\[
[ f ^ {*} (\alpha) ] ^ {2} = f ^ {*} ([ \alpha ] ^ {2}).\tag{28}
\]

Let \(\alpha, \beta\) be two differential forms of degree 1 on X with values in L of class \(C^{r-1}\) . The exterior product \(\alpha \wedge \beta\) of \(\alpha\) and \(\beta\) (Differentiable and Analytic Manifolds, R, 7.8.2) is a differential form of degree 2 on X with values in L of class \(C^{r-1}\) ; we have

\[
(\alpha \wedge \beta) _ {x} (u _ {1}, u _ {2}) = [ \alpha_ {x} (u _ {1}), \beta_ {x} (u _ {2}) ] - [ \alpha_ {x} (u _ {2}), \beta_ {x} (u _ {1}) ]\tag{29}
\]

for \(u_{1}, u_{2}\) in \(\mathrm{T}_x(\mathbf{X})\) . It is immediate that

\begin{enumerate}
\def\labelenumi{(\arabic{enumi})}
\setcounter{enumi}{29}
\tightlist
\item
\end{enumerate}

\[
[ \alpha + \beta ] ^ {2} = [ \alpha ] ^ {2} + [ \beta ] ^ {2} + \alpha \wedge \beta\tag{31}
\]

\[
\alpha \wedge \alpha = 2 [ \alpha ] ^ {2}.
\]

PROPOSITION 51. Let G be a Lie group, of finite dimension if K is of characteristic \textgreater0, and let \(a_{1},\ldots,a_{p}\) be elements of L(G), F a complete normable space and \(\alpha\) a differential form of degree \(p-1\) on G with values in F. If \(\alpha\) is left invariant, then

\[
(d \alpha) _ {e} (a _ {1}, \dots , a _ {p}) = \sum_ {i <   j} (- 1) ^ {i + j} \alpha_ {e} ([ a _ {i}, a _ {j} ], a _ {1}, \dots , a _ {i - 1}, a _ {i + 1}, \dots , a _ {j - 1},
\]

\[
\left. a _ {j + 1}, \dots , a _ {p}\right).
\]

If \(\alpha\) is right invariant, then

\[
\begin{array}{l} (d \alpha) _ {e} (a _ {1}, \dots , a _ {p}) \\ = - \sum_ {i <   j} (- 1) ^ {i + j} \alpha_ {e} ([ a _ {i}, a _ {j} ], a _ {1}, \dots , a _ {i - 1}, a _ {i + 1}, \dots , a _ {j - 1}, a _ {j + 1}, \dots , a _ {p}). \end{array}
\]

Suppose that \(\alpha\) is left invariant. By Differentiable and Analytic Manifolds, R, 8.5.7, then

\[
\begin{array}{l} (d \alpha) (L _ {a _ {1}}, \ldots , L _ {a _ {p}}) = \sum_ {i} (- 1) ^ {i - 1} L _ {a _ {i}} \alpha (L _ {a _ {1}}, \ldots , L _ {a _ {i - 1}}, L _ {a _ {i + 1}}, \ldots , L _ {a _ {p}}) \\ + \sum_ {i <   j} (- 1) ^ {i + j} \alpha ([ L _ {a _ {i}}, L _ {a _ {j}} ], L _ {a _ {1}}, \ldots , L _ {a _ {i - 1}}, L _ {a _ {i + 1}}, \ldots , L _ {a _ {j - 1}}, L _ {a _ {j + 1}}, \ldots , L _ {a _ {p}}). \end{array}
\]

But the functions \(\alpha (\mathrm{L}_{a_1},\dots ,\mathrm{L}_{a_{i - 1}},\mathrm{L}_{a_{i + 1}},\dots ,\mathrm{L}_{a_p})\) on G are left invariant and therefore constant. Hence

\[
\mathrm{L} _ {a _ {i}} \alpha (\mathrm{L} _ {a _ {1}}, \dots , \mathrm{L} _ {a _ {i - 1}}, \mathrm{L} _ {a _ {i + 1}}, \dots , \mathrm{L} _ {a _ {p}}) = 0.
\]

Moreover, \([L_{a_{i}}, L_{a_{j}}] = L_{\{a_{i}, a_{j}\}}\) (Proposition 23), whence the first formula of Proposition 51. The second can be established analogously, this time using the relation \([R_{a_{i}}, R_{a_{j}}] = -R_{\{a_{i}, a_{j}\}}\) .

COROLLARY 1. Let G be a Lie group, of finite dimension if K is of characteristic \textgreater0, and ω and ω' the left and right canonical differential forms of G. Then

\[
d \omega + [ \omega ] ^ {2} = 0, \quad d \omega^ {\prime} - [ \omega^ {\prime} ] ^ {2} = 0.
\]

By Proposition 51,

\[
\begin{array}{r l} (d \omega) _ {e} (a _ {1}, a _ {2}) & = - \omega_ {e} ([ a _ {1}, a _ {2} ]) = - [ a _ {1}, a _ {2} ] = - [ \omega_ {e} (a _ {1}), \omega_ {e} (a _ {2}) ] \\ & = - [ \omega ] _ {e} ^ {2} (a _ {1}, a _ {2}) \end{array}
\]

whence the first formula. The second can be established analogously.

COROLLARY 2. Suppose that G is finite-dimensional. Let \((e_1, \ldots, e_n)\) be a basis of L(G), \((e_1^*, \ldots, e_n^*)\) the dual basis, \((c_{ijk})\) the constants of structure of L(G) relative to the basis \((e_1, \ldots, e_n)\) and \(\omega_i\) (resp. \(\omega_i^{\prime})\) the left (resp. right) invariant differential form on G with values in K such that \((\omega_i)_e = e_i^*\) (resp. \((\omega_i^{\prime})_e = e_i^*)\) . Then

\[
d \omega_ {k} + \sum_ {i <   j} c _ {i j k} \omega_ {i} \wedge \omega_ {j} = 0 (k = 1, 2, \dots , n)
\]

\[
d \omega_ {k} ^ {\prime} - \sum_ {i <   j} c _ {i j k} \omega_ {i} ^ {\prime} \wedge \omega_ {j} ^ {\prime} = 0 (k = 1, 2, \dots , n).
\]

If \(r < s\) ,

\[
\begin{array}{r l} \left(d \omega_ {k}\right) _ {e} (e _ {r}, e _ {s}) & = - \left(\omega_ {k}\right) _ {e} ([ e _ {r}, e _ {s} ]) \\ & = - \sum_ {i} c _ {r s i} (\omega_ {k}) _ {e} (e _ {i}) \\ & = - c _ {r s k} \\ & = - \sum_ {i <   j} c _ {i j k} (\omega_ {i} \wedge \omega_ {j}) _ {e} (e _ {r}, e _ {s}). \end{array}
\]

The argument is similar for the \(\omega_{k}^{\prime}\) .

\subsection{CONSTRUCTION OF INVARIANT DIFFERENTIAL FORMS}

Lemma 2. Let \(\mathbf{G}\) be a Lie group, \(\mathrm{U}\) a symmetric open neighbourhood of \(e\) in \(\mathbf{G}\) , \(\mathrm{E}\) a complete normable space and \(\phi: \mathrm{U}^2 \to \mathrm{E}\) an analytic mapping. For all \(g \in \mathrm{U}\) , let \(\omega_g\) be the differential at the point \(g\) of the mapping \(h \mapsto \phi(g^{-1}h)\) . Then \(\omega\) is the restriction to \(\mathrm{U}\) of the left invariant differential form on \(\mathbf{G}\) whose value at \(e\) is \(d_e\phi\) .

Clearly \(\omega_{e} = d_{e}\phi\) . For all \(g \in \mathrm{U}\) and all \(t \in \mathrm{T}_{e}(\mathrm{G})\) ,

\[
\begin{array}{r l} \langle \omega_ {g}, \mathrm{T} _ {e} (\gamma (g)) t \rangle & = \langle d _ {g} (\phi \circ \gamma (g) ^ {- 1}), \mathrm{T} _ {e} (\gamma (g)) t \rangle \\ & = \langle d _ {e} \phi \circ \mathrm{T} _ {g} (\gamma (g) ^ {- 1}), \mathrm{T} _ {e} (\gamma (g)) t \rangle = \langle d _ {e} \phi , t \rangle \end{array}
\]

and hence \(\omega_{g}\) is derived from \(d_{e}\phi\) by \(\mathrm{T}_{e}(\gamma(g))\) .

PROPOSITION 52. Let \(n\) be an integer \(>0\) , \(G\) an \(n\) -dimensional Lie group, \(U\) a symmetric open neighbourhood of \(e\) in \(G\) and \(\psi: U^2 \to K^n\) a chart of \(G\) such that \(\psi(e) = 0\) . If \((x_1, \ldots, x_n)\) are the coordinates of \(x \in \psi(U)\) and \((y_1, \ldots, y_n)\) the coordinates of \(y \in \psi(U)\) , we denote by

\[
m _ {1} (x _ {1}, \dots , x _ {n}, y _ {1}, \dots , y _ {n}), \dots , m _ {n} (x _ {1}, \dots , x _ {n}, y _ {1}, \dots , y _ {n})
\]

the coordinates of \(\psi (\psi^{-1}(x)^{-1}\psi^{-1}(y))\) . Then, if we write, for \(1\leqslant k\leqslant n\)

\[
\begin{array}{r l} \varpi_ {k} (x _ {1}, \ldots , x _ {n}) = D _ {n + 1} m _ {k} (x _ {1}, \ldots , x _ {n}, x _ {1}, \ldots , x _ {n}) d x _ {1} + \dots \\ & \quad + \mathrm{D} _ {2 n} m _ {k} (x _ {1}, \ldots , x _ {n}, x _ {1}, \ldots , x _ {n}) d x _ {n}, \end{array}\tag{32}
\]

the differential forms \(\varpi_k\) on \(\psi(\mathbf{U})\) are derived through \(\psi\) from left invariant differential forms on \(\mathbf{G}\) and are such that \(\varpi_k(0,\ldots,0) = dx_k\) .

We apply Lemma 2 with E = K, taking \(\phi(g)\) to be the coordinate of \(\psi(g)\) of index k. We obtain a differential form \(\omega_{k}\) ; let \(\varpi_{k}\) be its transform under \(\psi\) . The value of \(\varpi_{k}\) at \((x_{1},\ldots,x_{n})\) is the differential at \((x_{1},\ldots,x_{n})\) of the function \(y\mapsto m_{k}(x_{1},\ldots,x_{n},y_{1},\ldots,y_{n})\) ; this value is therefore given by formula (32). It then suffices to use the conclusion to Lemma 2.

PROPOSITION 53. Let G be a Lie group, A a complete normable algebra and \(\phi\) a Lie group morphism of G into A*. For all \(g \in G\) , let \(\omega_g = \phi(g)^{-1}.d_g\phi\) . Then \(\omega\) is the left invariant differential form on G whose value at \(e\) is \(d_e\phi\) .

We apply Lemma 2 with \(\mathbf{E} = \mathbf{A}\) and \(\mathbf{U} = \mathbf{G}\) . The differential at \(g\) of the mapping \(h \mapsto \phi(g^{-1}h) = \phi(g)^{-1}\phi(h)\) is \(\phi(g)^{-1}.d_g\phi\) .

\subsection{HAAR MEASURE ON A LIE GROUP}

Let G be a Lie group of finite dimension n.~Then \(\bigwedge^{n}(\mathrm{T}_{e}(\mathrm{G}))\) is of dimension 1. Hence (no. 13) the vector space S of left invariant differential forms of degree n on G is of dimension 1. Let \((\omega_{1},\ldots,\omega_{n})\) be a basis of the space of left invariant differential forms of degree 1 on G; then \(\omega_{1}\wedge\omega_{2}\wedge\cdots\wedge\omega_{n}\) is a basis of S.

PROPOSITION 54. Let G be a Lie group of finite dimension n, \(\omega\) a left invariant differential form of degree n on G and \(\phi\) an endomorphism of G. Then

\[
\phi^ {*} (\omega) = (\det L (\phi)) \omega .
\]

We write \(\mathrm{L}(\phi) = u, \omega_e = f\) and \(\phi^*(\omega)_e = g\) . For all \(x_1, \ldots, x_n\) in \(\mathrm{L}(\mathrm{G})\) ,

\[
g (x _ {1}, \dots , x _ {n}) = f (u x _ {1}, \dots , u x _ {n}) = (\det u) f (x _ {1}, \dots , x _ {n})
\]

and hence \(\phi^{*}(\omega)_{e} = \det L(\phi). \omega_{e}\) . On the other hand, if \(g \in G\) ,

\[
\phi \circ \gamma (g) = \gamma (\phi (g)) \circ \phi
\]

and hence \(\gamma(g)^{*}\phi^{*}(\omega)=\phi^{*}(\omega)\) . Thus \(\phi^{*}(\omega)\) is left invariant, whence the proposition.

COROLLARY. For all \(g \in G\) ,

\[
\delta (g) ^ {*} \omega = (\det \operatorname{Ad} g) \omega .
\]

\[
\delta (g) ^ {*} \omega = \delta (g) ^ {*} \gamma (g) ^ {*} \omega = (\operatorname{Int} g) ^ {*} \omega \text {   and   } L (\operatorname{Int} g) = \operatorname{Ad} g.
\]

Let G be a locally compact group and \(\phi\) an endomorphism of G. Suppose that there exist open neighbourhoods V, \(V'\) of e such that \(\phi(V) = V'\) and \(\phi \mid V\) is a local isomorphism of G into G. Let \(\mu\) be a left Haar measure on G. By Integration, Chapter VII, §1, Corollary to Proposition 9, there exists a unique number a \textgreater{} 0 such that \(\phi(\mu \mid V) = a^{-1}\mu \mid V'\) . Clearly a is independent of the choice of V, \(V'\) and \(\mu\) . It is called the modulus of \(\phi\) and is denoted by mod \(_{G}\) \(\phi\) or simply mod \(\phi\) . When \(\phi\) is an automorphism of G, we recover Definition 4 of Integration, Chapter VII, §1.

PROPOSITION 55. Suppose that K is locally compact. Let \(\mu\) be a Haar measure on the additive group of K. Let G be a Lie group of finite dimension \(n\) .

\begin{enumerate}
\def\labelenumi{(\roman{enumi})}
\item
  Let \(\omega\) be a non-zero left invariant differential form of degree \(n\) on G. Then the measure mod \((\omega)_{\mu}\) (Differentiable and Analytic Manifolds, R, 10.1.6) is a left Haar measure on G. If K = R and G has the orientation defined by \(\omega\) , the measure defined by \(\omega\) (Differentiable and Analytic Manifolds, R, 10.4.3) is a left Haar measure on G.
\item
  Let \(\phi\) be an étale endomorphism of G. Then mod \(\phi = \text{mod}\det L(\phi)\) .
\item
  is obvious. Let \(\mathrm{V}, \mathrm{V}'\) be open neighbourhoods of \(e\) such that \(\phi(\mathrm{V}) = \mathrm{V}'\) and \(\phi \mid \mathrm{V}\) is a local isomorphism of \(\mathrm{G}\) into \(\mathrm{G}\) . Then
\end{enumerate}

\[
\begin{array}{r l r l} \phi^ {- 1} (\mathrm{mod} (\omega) _ {\mu}   |   \mathrm{V} ^ {\prime}) & = \mathrm{mod} (\phi^ {*} (\omega)) _ {\mu}   |   \mathrm{V}) & & \text {by transport of structure} \\ & = \mathrm{mod} (\det \mathrm{L} (\phi) \omega   |   \mathrm{V}) _ {\mu} & & \text {(Proposition 54)} \\ & = \mathrm{mod} \det \mathrm{L} (\phi) (\mathrm{mod} (\omega) _ {\mu}   |   \mathrm{V}) \end{array}
\]

whence mod \(\phi = \mathrm{mod}\det \mathrm{L}(\phi)\) by definition of mod \(\phi\) .

COROLLARY. For all \(g \in \mathbf{G}\) , \(\Delta_{\mathbf{G}}(g) = (\text{mod det } \operatorname{Ad} g)^{-1}\) . In particular, for \(\mathbf{G}\) to be unimodular, it is necessary and sufficient that mod det \(\operatorname{Ad} g = 1\) for all \(g \in \mathbf{G}\) .

\[
\begin{array}{r l} \Delta_ {\mathbf {G}} (g) & = (\text {mod} \operatorname{Int} g) ^ {- 1} \quad \text {(Integration, Chapter VII, § 1, formula (33))} \\ & = (\text {mod} \det \operatorname{L} (\operatorname{Int} g)) ^ {- 1} \quad \text {(Proposition 55)} \\ & = (\text {mod} \det \operatorname{Ad} g) ^ {- 1}. \end{array}
\]

Remark. Preserving the hypotheses and notation of Proposition 52, suppose that K is locally compact. Let \(\mu\) be the measure

\[
\operatorname{mod} \det (\mathrm{D} _ {n + i} m _ {k} (x _ {1}, \dots , x _ {n}, x _ {1}, \dots , x _ {n})) _ {1 \leqslant i, k \leqslant n} d x _ {1} \dots d x _ {n}
\]

on \(\psi (\mathbf{U})\) . Then \(\psi^{-1}(\mu)\) is the restriction to U of a Haar measure on G.

PROPOSITION 56. Let G be a Lie group of finite dimension n, H a p-dimensional Lie subgroup and X the Lie homogeneous space G/H. Suppose that

\[
\det \operatorname{Ad} _ {\mathrm{L} (\mathrm{G})} h = \det \operatorname{Ad} _ {\mathrm{L} (\mathrm{H})} h
\]

for all \(h\in \mathbf{H}\) . Then:

\begin{enumerate}
\def\labelenumi{(\roman{enumi})}
\item
  The differential forms of degree \(n - p\) on \(\mathbf{X}\) which are invariant under \(\mathbf{G}\) are analytic.
\item
  The vector space of these forms is of dimension 1.
\item
  If \(\omega\) is such a non-zero form and K is locally compact, mod \((\omega)_{\mu}\) is a non-zero measure on X which is invariant under G.
\end{enumerate}

By § 1, no. 8, Examples, \(\mathrm{Alt}^{n - p}(\mathrm{TX},\mathrm{K})\) is an analytic vector G-bundle. Let \(x_0\) be the canonical image of \(e\) in \(\mathbf{X}\) ; its stabilizer is H. The fibre of \(\mathrm{Alt}^{n - p}(\mathrm{TX},\mathrm{K})\) at \(x_0\) is \(\bigwedge^{n - p}\mathrm{T}_{x_0}(\mathbf{X})^*\) and \(\mathrm{T}_{x_0}(\mathbf{X})\) is canonically identified with \(\mathrm{L}(G) / L(\mathrm{H})\) . If \(h\in \mathrm{H}\) , the automorphism \(\tau_h\) of \(\mathbf{X}\) defined by \(h\) is derived when passing to the quotient from the automorphism \(g\mapsto hgh^{-1}\) of G. Hence the automorphism \(\mathrm{T}_{x_0}(\tau_h)\) is derived when passing to the quotient from \(\mathrm{Ad}_{\mathrm{L(G)}}(h)\) . As

\[
\det \mathrm{Ad} _ {\mathrm{L(G)}} h = (\det \mathrm{Ad} _ {\mathrm{L(H)}} h). (\det \mathrm{T} _ {x _ {0}} (\tau_ {h})),
\]

the hypothesis implies that \(\det \mathrm{T}_{x_0}(\tau_h) = 1\) . Thus, every element of \(\bigwedge^{n - p}\mathrm{T}_{x_0}(\mathbf{X})^*\) is invariant under H. Then (i) and (ii) follow from § 1, no. 8, Corollary 1 to Proposition 17 and (iii) is obvious.

The existence of a non-zero positive measure on X invariant under G follows from Integration, Chapter VII, § 2, Corollary 2 to Theorem 3, for the hypothesis of Proposition 56 implies \(\Delta_{\mathbf{G}}|H = \Delta_{\mathbf{H}}\) (Corollary to Proposition 55).

PROPOSITION 57. Let G be a Lie group of finite dimension n.~Choose a basis for \(\bigwedge^{n}\mathrm{T}_{e}(\mathrm{G})^{*}\) ; by means of the right (resp. left) trivialization of \(\bigwedge^{n}\mathrm{T}(\mathrm{G})^{*}\) , we can identify this vector bundle with the trivial vector bundle \(\mathbf{G} \times \mathbf{K}\) , so that the transpose of a scalar differential operator is identified with a scalar differential operator.

Then, if \(u \in \mathrm{U}(\mathrm{G})\) , the transpose of \(\mathrm{L}_u\) (resp. \(\mathbf{R}_u\) ) is \(\mathrm{L}_u^\vee\) (resp. \(\mathbf{R}_u^\vee\) ).

We shall consider the case where \(\bigwedge^{n}\mathrm{T}(\mathrm{G})^{*}\) has been trivialized using a right invariant form \(\omega\) .

Suppose that the proposition has been proved for elements \(u_{1}, u_{2}\) of U(G). Then,

\[
\begin{array}{l l} ^ {t} (\mathrm{L} _ {u _ {1} * u _ {2}}) = ^ {t} (\mathrm{L} _ {u _ {1}} \circ \mathrm{L} _ {u _ {2}}) & \text {(Proposition 23)} \\ = ^ {t} (\mathrm{L} _ {u _ {2}}) \circ^ {t} (\mathrm{L} _ {u _ {1}}) & \text {(Diff. \& Anal. Man., R, 14.3.3)} \\ = \mathrm{L} _ {u _ {2}} ^ {\vee} \circ \mathrm{L} _ {u _ {1}} ^ {\vee} & \text {by hypothesis} \\ = \mathrm{L} _ {u _ {2} * u _ {1}} ^ {\vee \vee} & \text {(Proposition 23)} \\ = \mathrm{L} _ {(u _ {1} * u _ {2}) \vee} & \text {(Proposition 7)} \end{array}
\]

and hence the proposition is true for \(u_{1} * u_{2}\) . It therefore suffices to prove the proposition when \(u \in \mathrm{T}_{e}(\mathrm{G})\) . Now \(\mathbf{L}_u\) is defined by G operating on G on the right (no. 6) and hence \(\theta_{\mathrm{L}_u}\omega = 0\) since \(\omega\) is right invariant (Differentiable and Analytic Manifolds, R, 8.4.5); therefore, if \(f\) is an analytic function in an open neighbourhood of \(e\) with values in K, then \(\theta_{\mathrm{L}_u}(f\omega) = (\theta_{\mathrm{L}_u}f)\omega\) (Differentiable and Analytic Manifolds, R, 8.4.8). Using the identifications made and Differentiable and Analytic Manifolds, R, 14.4.1, the transpose of \(\mathbf{L}_u\) is \(-\mathbf{L}_u\) , that is \(\mathbf{L}_u^\vee\) .

COROLLARY. Let G be a finite-dimensional real Lie group, \(\mu\) (resp. \(\nu\)) a left (resp. right) Haar measure on G, k an integer \(\geqslant0\) , \(u\in\mathrm{U}_{k}(\mathrm{G})\) and f and g real-valued functions of class \(C^{k}\) on G with compact support. Then

\[
\int_ {G} \left(R _ {u} f\right) g d \mu = \int_ {G} f (R _ {u} ^ {\vee} g) d \mu
\]

\[
\int_ {\mathbf {G}} \left(\mathrm{L} _ {u} f\right) g d \nu = \int_ {\mathbf {G}} f (\mathrm{L} _ {u} ^ {\vee} g) d \nu .
\]

This follows from Proposition 57 and Differentiable and Analytic Manifolds, R, 14.3.8.

\subsection{LEFT DIFFERENTIAL}

DEFINITION 9. Let G be a Lie group, M a manifold of class \(\mathbf{C}^r\) and \(f\) a mapping of class \(\mathbf{C}^r\) of M into G. The left (resp. right) differential of \(f\) is the differential form of degree 1 on M with values in L(G) which associates with every vector \(u \in \mathrm{T}_m(\mathrm{M})\) the element \(f(m)^{-1}\) . \((\mathrm{T}_m f)(u)\) (resp. \((\mathrm{T}_m f)(u). f(m)^{-1}\) ).

In this chapter we shall only consider the left differential, which we shall denote by \(f^{-1}.df\) , and leave to the reader the task of translating the results for the right differential.

If \(f\) is the identity mapping of \(G, f^{-1}.df\) is the canonical left differential form \(\omega\) of \(G\) . Returning to the general case of Definition 8,

\[
(f ^ {- 1}. d f) _ {m} = \omega_ {f (m)} \circ \mathrm{T} _ {m} (f)
\]

and hence \(f^{-1}.df = f^{*}(\omega)\) . This implies that \(f^{-1}.df\) is of class \(\mathbf{C}^{r - 1}\) .

Examples. (1) If G is the additive group of a complete normable space and \(T_{0}(E)\) is canonically identified with E, \(f^{-1}\) . df is the differential df defined in Differentiable and Analytic Manifolds, R, 8.2.2.

\begin{enumerate}
\def\labelenumi{(\arabic{enumi})}
\setcounter{enumi}{1}
\tightlist
\item
  Suppose that G is the multiplicative group A* associated with a complete normable algebra A. Then f can be considered as a mapping of M into A and hence the differential df in the sense of Differentiable and Analytic Manifolds, R, 8.2.2 is defined and the product \(f^{-1}df\) in the sense of Differentiable and Analytic Manifolds, R, 8.3.2 is defined. Clearly the latter form is identical with the left differential of f.
\end{enumerate}

PROPOSITION 58. Let G and H be two Lie groups, M a manifold of class \(\mathbf{C}^r\) , \(f\) a mapping of class \(\mathbf{C}^r\) of M into G and \(h\) a morphism of G into H. Then

\[
(h \circ f) ^ {- 1}. d (h \circ f) = \mathrm{L} (h) \circ (f ^ {- 1}. d f) = (h ^ {- 1}. d h) \circ \mathrm{T} (f).
\]

For all \(x \in \mathbf{M}\) and \(u \in \mathrm{T}_x(\mathrm{M})\) ,

\[
(h \circ f) ^ {- 1}. d (h \circ f) (u) = ((h \circ f) (x)) ^ {- 1}. T (h \circ f) (u).
\]

The latter expression is equal, on the one hand, to

\[
\begin{array}{l l} \mathrm{T} (h) (f (x) ^ {- 1}. \mathrm{T} (f) (u)) & (\S 2, \text { Proposition   5 }) \\ = \mathrm{T} _ {e} (h) ((f ^ {- 1}. d f) (u)) \end{array}
\]

and, on the other hand, to

\[
\begin{array}{c} h (f (x)) ^ {- 1} \mathrm{T} (h) (\mathrm{T} (f) u) \\ = (h ^ {- 1}. d h) (\mathrm{T} (f) u). \end{array}
\]

PROPOSITION 59. Let G be a Lie group, M a manifold of class \(\mathbf{C}^r\) , \(f\) and \(g\) mappings of class \(\mathbf{C}^r\) of M into G and \(p\) the canonical surjection of TM onto M.

\[
(i) (f g) ^ {- 1}. d (f g) = (\mathrm{Ad} \circ g \circ p) ^ {- 1} \circ (f ^ {- 1}. d f) + g ^ {- 1}. d g.
\]

\begin{enumerate}
\def\labelenumi{(\roman{enumi})}
\setcounter{enumi}{1}
\tightlist
\item
  Writing \(h(m) = f(m)^{-1}\) for all \(m \in M\) ,
\end{enumerate}

\[
h ^ {- 1}. d h = - (\mathrm{Ad} \circ f \circ p) \circ (f ^ {- 1}. d f).
\]

Assertion (i) follows from § 2, no. 2, Proposition 7. Assertion (ii) follows from (i) by putting \(g = h\) .

COROLLARY 1. Let \(s \in G\) and \(sg\) be the mapping \(x \mapsto sg(x)\) of \(M\) into \(G\) . Then \((sg)^{-1}.d(sg) = g^{-1}.dg\) .

This follows from Proposition 59 (i) taking \(f\) to be the constant mapping \(x \mapsto s\) of \(M\) into \(G\) .

COROLLARY 2. If the mappings \(f\) and \(g\) of \(M\) into \(G\) have the same left differential, the tangent mapping to \(fg^{-1}\) is everywhere zero. If further \(K\) is of characteristic 0, then \(fg^{-1}\) is locally constant.

By Proposition 59,

\[
(f g ^ {- 1}) ^ {- 1}. d (f g ^ {- 1}) = (\mathrm{Ad} \circ g \circ p) \circ (f ^ {- 1}. d f) - (\mathrm{Ad} \circ g \circ p) \circ (g ^ {- 1}. d g).
\]

If \(f^{-1}.df = g^{-1}.dg\) , then \((fg^{-1})^{-1}.d(fg^{-1}) = 0\) , that is \(T_x(fg^{-1}) = 0\) for all \(x \in M\) . This proves the first assertion. The second follows from it by Differentiable and Analytic Manifolds, R, 5.5.3.

PROPOSITION 60. Let G be a Lie group, of finite dimension if K is of characteristic \textgreater0, M a manifold of class \(\mathbf{C}^r\) , \(f\) a mapping of class \(\mathbf{C}^r\) of M into G and \(\alpha\) the left differential of \(f\) . Then \(d\alpha + [\alpha]^2 = 0\) .

Let \(\omega\) be the canonical left differential form of G. Using Corollary 1 to Proposition 51, no. 14, we have

\[
\begin{array}{r l} d \alpha & = d (f ^ {*} (\omega)) = f ^ {*} (d \omega) = f ^ {*} (- [ \omega ] ^ {2}) \\ & = - [ f ^ {*} (\omega) ] ^ {2} = - [ \alpha ] ^ {2}. \end{array}
\]

\subsection{LIE ALGEBRA OF A LIE GROUP GERM}

In this no. \((G, e, \theta, m)\) denotes a Lie group germ. A large part of the results of the § are still true with the same proof. We shall review those which we shall find useful.

18.1. Let \(\Omega\) be the set of definition of \(m\) . Let \((g, g') \in \Omega\) , \(t \in \mathrm{T}_g^{(\alpha)}(\mathrm{G})\) , \(t' \in \mathrm{T}_{g'}^{(\alpha)}(\mathrm{G})\) . As in no. 1, the convolution product of \(t\) and \(t'\) , denoted by \(t * t'\) , is the image of \(t \otimes t'\) under \(m\) . We write \(\mathrm{U}(\mathrm{G}) = \mathrm{T}_e^{(\alpha)}(\mathrm{G})\) , \(\mathrm{U}_s(\mathrm{G}) = \mathrm{T}_e^{(s)}(\mathrm{G})\) , \(\mathrm{U}^+(\mathrm{G}) = \mathrm{T}_e^{(\alpha)+}(\mathrm{G})\) , \(\mathrm{U}_s^+(\mathrm{G}) = \mathrm{T}_e^{(s)+}(\mathrm{G})\) . For \(t\) , \(t'\) in \(\mathrm{U}(\mathrm{G})\) , \(t * t'\) is defined and belongs to \(\mathrm{U}(\mathrm{G})\) . With the convolution product, \(\mathrm{U}(\mathrm{G})\) is an associative algebra with unit element \(\varepsilon_e\) , filtered by the \(\mathrm{U}_s(\mathrm{G})\) . The canonical isomorphism \(i_{\mathrm{G}, e}\) of gr \(\mathrm{U}(\mathrm{G})\) onto \(\mathrm{TS}(\mathrm{T}_e(\mathrm{G}))\) is an algebra isomorphism.

18.2. Let G, H be Lie group germs and \(\phi: G \to H\) a morphism. If \(t \in \mathrm{U}(\mathrm{G})\) , the image \(\mathrm{U}(\phi)(t)\) of t under \(\phi_{*}\) is an element of U(H) and \(\mathrm{U}(\phi)\) is a morphism of the algebra U(G) into the algebra U(H). The mapping \(\theta: x \mapsto x^{-1}\) of G into G defines a mapping \(t \mapsto t^{\vee}\) of U(G) into U(G). For \(t, t'\) in U(G), the product \(t * t'\) evaluated relative to \(G^{\vee}\) is equal to the product \(t' * t\) evaluated relative to G and \((t * t')^{\vee} = t^{\prime\vee} * t^\vee\) . Then \(\mathrm{U}(\phi)(t^{\vee}) = (\mathrm{U}(\phi)t)^{\vee}\) . If \(G_{1}, \ldots, G_{n}\) are Lie group germs and \(G = G_{1} \times \cdots \times G_{n}\) , the canonical isomorphism of \(\mathrm{U}(G_{1}) \otimes \cdots \otimes \mathrm{U}(G_{n})\) onto U(G) is an algebra isomorphism;

for \(t_{1},\ldots,t_{n}\) in U(G), \((t_{1}\otimes\cdots\otimes t_{n})^{\vee}=t_{1}^{\vee}\otimes\cdots\otimes t_{n}^{\vee}\) . Let H be a Lie subgroup germ of G and i: H→G the canonical injection. Then U(i) is an injective homomorphism of the algebra U(H) into U(G) and

\[
\mathrm{U} (i) \left(t ^ {\vee}\right) = \left(\mathrm{U} (i) (t)\right) ^ {\vee}
\]

for all \(t \in \mathrm{U}(\mathrm{H})\) . With the convolution product and the coproduct defined by the manifold structure on G, \(\mathrm{U}(\mathrm{G})\) is a bigebra and \(\mathrm{U}(\phi)\) is a bigebra morphism.

18.3. Let G be a Lie group germ, X a manifold of class \(C^{r}\) and \(\psi\) a law chunk of left operation of class \(C^{r}\) of G on X. Let \(\Omega\) be the set of definition of \(\psi\) . If \(t \in T_{g}^{(s)}(G)\) , \(u \in T_{x}^{(s^{\prime\prime})}(X)\) , \((g, x) \in \Omega\) and \(s + s^{\prime} \leqslant r\) , let \(t * u\) denote the image of \(t \otimes u\) under \(\psi_{*}\) . Let \(t \in T_{g}^{(s)}(G)\) , \(t^{\prime} \in T_{g'}^{(s^{\prime})}(G)\) , \(u \in T_{x}^{(s^{\prime})}(X)\) ; if \(s + s^{\prime} + s^{\prime\prime} \leqslant r\) and \(gg^{\prime}\) , \((gg^{\prime})x\) , \(g^{\prime}x\) , \(g(g^{\prime}x)\) are defined, then

\[
(t * t ^ {\prime}) * u = t * (t ^ {\prime} * u).
\]

Let \(x_0 \in \mathbf{X}\) and \(\rho(x_0)\) be the mapping \(g \mapsto gx_0\) , which is defined in an open neighbourhood of \(e\) . If \(t \in \mathrm{U}_r(\mathrm{G})\) , then \(\rho(x_0)_*t = t * \varepsilon_{x_0}\) . Here and in the rest of this no., we shall leave to the reader the task of translating the results for law chunks of right operation.

18.4. Preserving the notation of 18.3, let \(t \in \mathrm{U}_s(\mathbf{G})\) with \(s \leqslant r\) . Let \(f\) be a function of class \(\mathbf{C}^r\) on \(\mathbf{X}\) with values in a Hausdorff polynormed space. Let \(t * f\) denote the function on \(\mathbf{X}\) defined by

\[
\begin{array}{r l} (t * f) (x) & = \langle t, g \mapsto f (\psi (\theta (g), x)) \rangle \\ & = \langle t ^ {\vee}, f \circ \rho (x) \rangle = \langle \rho (x) _ {*} (t ^ {\vee}), f \rangle = \langle t ^ {\vee} * \varepsilon_ {x}, f \rangle . \end{array}
\]

If \(t \in \mathrm{U}_s(\mathrm{G})\) , \(t' \in \mathrm{U}_{s'}(\mathrm{G})\) and \(s + s' \leqslant r\) , then \(\langle t', t * f \rangle = \langle t^\vee * t', f \rangle\) and \((t * t') * f = t * (t' * f)\) . Let \(t \in \mathrm{U}_s(\mathrm{G})\) , \(f\) and \(f'\) be functions of class \(C^r\) on X with values in Hausdorff polynormed spaces \(F, F'\) and \((u, u') \mapsto u. u'\) be a continuous bilinear mapping of \(F \times F'\) into a Hausdorff polynormed space; let \(\underline{n}\) .

\(\sum_{i=1}^{n} t_i \otimes t_i'\) be the image of \(t\) under the coproduct; if \(s \leqslant r\) , then

\[
t * (f f ^ {\prime}) = \sum_ {i = 1} ^ {n} (t _ {i} * f) (t _ {i} ^ {\prime} * f ^ {\prime}).
\]

18.5. Preserving the notation of 18.3, let \(t \in \mathrm{U}_s(\mathrm{G})\) with \(s \leqslant r\) . The mapping \(x \to t * \varepsilon_x\) is called the field of point distributions defined by \(t\) and the law chunk of operation and is sometimes denoted by \(\mathrm{D}_t^{\psi}\) or \(\mathrm{D}_t\) . If \(f: \mathrm{X} \to \mathrm{F}\) is a function of class \(C^r\) , the function \(t^{\vee} * f\) on \(\mathrm{X}\) is also denoted by \(\mathrm{D}_t f\); it is of class \(C^{r - s}\) if \(s < \infty\) . If \(t \in \mathrm{U}_s(\mathrm{G})\) , \(t' \in \mathrm{U}_{s'}(\mathrm{G})\) and \(s + s' \leqslant r\) , then \(\mathrm{D}_{t,t'} f = \mathrm{D}_{t'}(\mathrm{D}_t f)\) . If \(G\) and \(X\) are finite-dimensional, \(D_t\) is a differential operator on \(X\) of order \(\leqslant s\) and of class \(C^{r - s}\) (if \(s < \infty\) ). The function \(D_t f\) is then the transform of \(f\) under this differential operator.

18.6. Let G be a Lie group germ and \(t \in \mathrm{U}(\mathrm{G})\) . \(L_{t}\) denotes the field of point distributions \(g \mapsto \varepsilon_{g} * t\) on G and \(R_{t}\) the field of point distributions \(g \mapsto t * \varepsilon_{g}\) on G. If \(f \in \mathcal{C}^{\omega}(\mathrm{G}, \mathrm{F})\) , then \(\mathrm{L}_{t} f \in \mathcal{C}^{\omega}(\mathrm{G}, \mathrm{F})\) and \(\mathrm{R}_{t} f \in \mathcal{C}^{\omega}(\mathrm{G}, \mathrm{F})\) . For \(t, t'\) in \(\mathrm{U}(\mathrm{G})\) , \(L_{t* t'} = L_t \circ L_{t'}, R_{t* t'} = R_{t'} \circ R_t, L_t \circ R_{t'} = R_{t'} \circ L_t, \theta(L_t) = R_t^\vee\) .

18.7. As \(T_{e}(G)\) is the set of primitive elements of U(G),

\[
[ \mathrm{T} _ {e} (\mathrm{G}), \mathrm{T} _ {e} (\mathrm{G}) ] \subset \mathrm{T} _ {e} (\mathrm{G}).
\]

The normable space \(\mathrm{T}_{e}(\mathrm{G})\) , together with the bracket, is a normable Lie algebra, called the normable Lie algebra of G (or Lie algebra of G) and denoted by \(\mathrm{L}(\mathrm{G})\) . Let \(\mathrm{E}(\mathrm{G})\) be the enveloping algebra of \(\mathrm{L}(\mathrm{G})\) . The canonical injection of \(\mathrm{L}(\mathrm{G})\) into \(\mathrm{U}(\mathrm{G})\) defines a homomorphism \(\eta\) of the algebra \(\mathrm{E}(\mathrm{G})\) into the algebra \(\mathrm{U}(\mathrm{G})\) ; if K is of characteristic 0, \(\eta\) is a bigebra isomorphism, by means of which \(\mathrm{U}(\mathrm{G})\) is identified with \(\mathrm{E}(\mathrm{G})\) . Using the notation of 18.3, for all \(a \in \mathrm{L}(\mathrm{G})\) , let \(D_{a}\) be the field of point distributions defined by a on X. The mapping \((a, x) \mapsto D_{a}(x)\) is a morphism of class \(C^{r-1}\) of the trivial vector bundle \(\mathrm{L}(\mathrm{G}) \times \mathrm{X}\) into the vector bundle \(\mathrm{T}(\mathrm{X})\) . Let I be an open subset of K containing 0 and \(\gamma: I \to G\) a mapping of class \(C^{r}\) such that \(\gamma(0) = e\) . Let \(a = T_{0}(\gamma) 1 \in L(\mathrm{G})\) . If f: X → F is a function of class \(C^{r}\) , then

\[
\left(\mathrm{D} _ {a} f\right) (x) = \lim _ {k \in \mathrm{K} ^ {*}, k \rightarrow 0} k ^ {- 1} \left(f (\gamma (k) x) - f (x)\right).
\]

If \(r \geqslant 2\) , the mapping \(a \mapsto \mathrm{D}_a\) is a law of left infinitesimal operation of class \(\mathrm{C}^{r-1}\) of \(\mathrm{L}(\mathrm{G})\) on \(\mathrm{X}\) .

18.8. Let G and H be Lie group germs and \(\phi\) a morphism of G into H. The restriction of \(\mathrm{U}(\phi)\) to \(\mathrm{L}(\mathrm{G})\) , which is just \(\mathrm{T}_{e}(\phi)\) , is a continuous morphism of \(\mathrm{L}(\mathrm{G})\) into \(\mathrm{L}(\mathrm{H})\) , which we denote by \(\mathrm{L}(\phi)\) . If \(\psi\) is a morphism of H into a Lie group germ, then \(\mathrm{L}(\psi \circ \phi) = \mathrm{L}(\psi) \circ \mathrm{L}(\phi)\) . For \(\phi\) to be an immersion, it is necessary and sufficient that \(\mathrm{L}(\phi)\) be an isomorphism of \(\mathrm{L}(\mathrm{G})\) onto a Lie subalgebra of \(\mathrm{L}(\mathrm{H})\) which admits a topological supplement. In particular, if G is a Lie subgroup germ of H and \(\phi\) is the canonical injection, \(\mathrm{L}(\mathrm{G})\) is identified with a Lie subalgebra of \(\mathrm{L}(\mathrm{H})\) by means of \(\mathrm{L}(\phi)\) . If \((\mathrm{G}_{i})_{i\in I}\) is a finite family of Lie group germs and G is their product, \(\mathrm{L}(\mathrm{G})\) is canonically identified with \(\prod_{i\in I}\mathrm{L}(\mathrm{G}_{i})\) .

18.9. Let G be a Lie group germ, of finite dimension if K is of characteristic \textgreater0. Let F be a complete normable space. Let \(\alpha\) be a differential form of degree k on G with values in F. \(\alpha\) is called left invariant on G if \(\alpha_{g}\) is derived from \(\alpha_{e}\) by the mapping \(h \mapsto gh\) of a neighbourhood of e onto a neighbourhood of g. If \(\alpha\) is left invariant, \(\alpha\) is analytic. The mapping \(\alpha \mapsto \alpha_{e}\) is a bijection of the set of left invariant differential forms of degree k on G with values in F onto the set of continuous alternating k-linear mappings of \(T_{e}(G)\) into F. If \(\alpha_{e} = \operatorname{Id}_{\mathrm{T}_e(\mathrm{G})}\) , \(\alpha\) is called the left canonical differential form of G. The definitions of right invariant differential forms and the right canonical differential form of G are analogous. If \(\omega\) is the left canonical differential form of G, then \(d\omega + [\omega]^{2} = 0\) . Let M be a manifold of class \(C^{r}\) and f a mapping of class \(C^{r}\) of M into G. The left differential of f, denoted by \(f^{-1}.df\) , is the differential form of degree 1 on M with values in L(G) which associates with each vector \(u \in T_{m}(M)\) the element \(f(m)^{-1}\) . \((\mathrm{T}_{m}f)(u)\) . Then \(f^{-1}.df = f^{*}(\omega)\) and \(d\alpha + [\alpha]^{2} = 0\) . If two mappings f and g of M into G have the same left differential and K is of characteristic 0, then \(fg^{-1}\) is locally constant.

\section{PASSAGE FROM LIE ALGEBRAS TO LIE GROUPS}

Recall that, until the end of the chapter, K is assumed to be of characteristic 0.

\subsection{PASSAGE FROM LIE ALGEBRA MORPHISMS TO LIE GROUP MORPHISMS}

Lemma 1. Let \(\mathbf{G}\) be a Lie group germ and \(\mathfrak{h}\) a Lie subalgebra of \(\mathbf{L}(\mathbf{G})\) admitting a topological supplement. The union of the \(\mathfrak{gh}\) (resp. \(\mathfrak{hg}\) ) for \(g\in\mathbf{G}\) is an integrable vector subbundle of \(\mathbf{T}(\mathbf{G})\) .

By considering the left trivialization of \(\mathrm{T}(\mathrm{G})\) (§ 2, no. 3), it is seen immediately that the \(g\mathfrak{h}\) , for \(g\in \mathrm{G}\) , are the fibres of a vector subbundle \(\mathrm{E}\) of \(\mathrm{T}(\mathrm{G})\) . Let \(g\in \mathrm{G}\) . The set of \((\mathrm{L}_a)_g\) , where \(a\in \mathfrak{h}\) , is equal to \(g\mathfrak{h}\) . Now, if \(a\) and \(b\) belong to \(\mathfrak{h}\) , then \([\mathrm{L}_a,\mathrm{L}_b] = \mathrm{L}_{[a,b)}\) and \([a,b]\in \mathfrak{h}\) . Hence \(\mathrm{E}\) is integrable (Differentiable and Analytic Manifolds, R, 9.3.3 (iv)). The argument is similar for the \(\mathfrak{h}g\) .

The integral foliation (Differentiable and Analytic Manifolds, R, 9.3.2) of the union of the \(\mathfrak{gh}\) (resp. \(\mathfrak{hg}\) ) is called the left (resp. right) foliation of G associated with \(\mathfrak{h}\) .

THEOREM 1. Let \(\mathbf{G}\) and \(\mathbf{H}\) be Lie group germs and \(f\) a continuous morphism of \(\mathbf{L}(\mathbf{G})\) into \(\mathbf{L}(\mathbf{H})\) .

\begin{enumerate}
\def\labelenumi{(\roman{enumi})}
\item
  There exist an open Lie subgroup germ \(\mathbf{G}'\) of \(\mathbf{G}\) and a morphism \(\phi\) of \(\mathbf{G}'\) into \(\mathbf{H}\) such that \(f = \mathbf{L}(\phi)\) .
\item
  Let \(\mathrm{G}_1, \mathrm{G}_2\) be open Lie subgroup germs of G and \(\phi_i\) a morphism of \(\mathrm{G}_i\) into H such that \(f = \mathrm{L}(\phi_i)\) for \(i = 1, 2\) . Then \(\phi_1\) and \(\phi_2\) coincide on a neighbourhood of \(e\) . Let \(p_1: \mathrm{G} \times \mathrm{H} \to \mathrm{G}\) , \(p_2: \mathrm{G} \times \mathrm{H} \to \mathrm{H}\) be the canonical projections. For all \((g,h) \in \mathrm{G} \times \mathrm{H}\) , let \(f_{g,h}\) be the mapping \(ga \mapsto hf(a)\) of \(\mathrm{T_g(G)} = g\mathrm{L(G)}\) into \(\mathrm{T_h(H)} = h\mathrm{L(H)}\) . By considering the left trivializations of \(\mathrm{T(G)}\) and \(\mathrm{T(H)}\) , it is seen immediately that the \(f_{g,h}\) define a morphism of \(p_1^*\mathrm{T(G)}\) into \(p_{2}^{*}\mathbf{T}(\mathrm{H})\) . Let \(a\) be the graph of \(f\) ; it is a closed Lie subalgebra of \(\mathbf{L}(\mathrm{G}) \times \mathbf{L}(\mathrm{H})\) which admits \(\{0\} \times \mathbf{L}(\mathrm{H})\) as topological supplement. For all \((g,h) \in \mathrm{G} \times \mathrm{H}\) , the graph of \(f_{g,h}\) is \((g,h)\) . \(a\) . The union of these graphs is an integrable vector subbundle of \(\mathrm{T}(\mathrm{G} \times \mathrm{H})\) (Lemma 1). Then there exist (Differentiable and Analytic Manifolds, R, 9.3.7) an open neighbourhood \(\mathrm{U}\) of \(e_{\mathrm{G}}\) in \(\mathrm{G}\) and an analytic mapping \(\phi\) of \(\mathrm{U}\) into \(\mathrm{H}\) such that \(\phi(e_{\mathrm{G}}) = e_{\mathrm{H}}\) and \(\mathrm{T}_{g}(\phi) = f_{g,\phi(g)}\) for all \(g \in \mathrm{U}\) . In particular, \(\mathrm{T}_{e_{\mathrm{G}}}(\phi) = f\) .
\end{enumerate}

Let \(\mathbf{V}\) be an open neighbourhood of \(e_{\mathrm{G}}\) in \(\mathbf{G}\) such that, for \((s,t)\in \mathbf{V}\times \mathbf{V}\) , the products \(st\) and \(\phi (s)\phi (t)\) are defined and \(st\in \mathbf{U}\) . Consider the mappings \(\alpha_{1},\alpha_{2}\) of \(\mathbf{V}\times \mathbf{V}\) into \(\mathbf{H}\) defined by

\[
\alpha_ {1} (s, t) = \phi (t s), \quad \alpha_ {2} (s, t) = \phi (t) \phi (s).
\]

Then \(\alpha_{1}(t,e) = \phi (t) = \alpha_{2}(t,e)\) . On the other hand, let \(t\) be fixed in V and \(\beta_{i}\) be the mapping \(s\mapsto \alpha_i(s,t)\) of V into H. Then, for all \(s\in \mathrm{V}\) and all \(a\in \mathrm{L}(\mathrm{G})\) ,

\[
\begin{array}{r l} \mathrm{T} _ {s} (\beta_ {1}) (s a) & = \mathrm{T} _ {t s} (\phi) (t s a) = f _ {t s, \phi (t s)} (t s a) \\ & = \phi (t s) f (a) = f _ {s, \beta_ {1} (s)} (s a) \\ \mathrm{T} _ {s} (\beta_ {2}) (s a) & = \phi (t) \mathrm{T} _ {s} (\phi) (s a) = \phi (t) f _ {s, \phi (s)} (s a) \\ & = \phi (t) \phi (s) f (a) = f _ {s, \beta_ {2} (s)} (s a). \end{array}
\]

Hence (Differentiable and Analytic Manifolds, R, 9.3.7) \(\alpha_{1}\) and \(\alpha_{2}\) coincide on a neighbourhood of \((e_{\mathrm{G}}, e_{\mathrm{G}})\) . The restriction of \(\phi\) to a sufficiently small symmetric open neighbourhood of \(e_{\mathrm{G}}\) is therefore a morphism of Lie group germs, whence (i).

Let \(\mathrm{G}_1, \mathrm{G}_2, \phi_1, \phi_2\) be as in (ii) and let us prove that \(\phi_1, \phi_2\) coincide on a neighbourhood of \(e_{\mathrm{G}}\) . There exists an open neighbourhood W of \(e_{\mathrm{G}}\) such that \(\phi_1(t s) = \phi_1(t)\phi_1(s)\) , \(\phi_2(t s) = \phi_2(t)\phi_2(s)\) for all \(s\) , \(t\) in W. Then, if \(s \in W\) and \(a \in L(G)\) ,

\[
\mathrm{T} _ {s} \left(\phi_ {i}\right) (s a) = \phi_ {i} (s) \mathrm{T} _ {e} \left(\phi_ {i}\right) (a) = \phi_ {i} (s) f (a) = f _ {s, \phi _ {i} (s)} (s a)
\]

for \(i = 1,2\) . As \(\phi_1(e_{\mathrm{G}}) = e_{\mathrm{H}} = \phi_2(e_{\mathrm{G}})\) , it follows from Differentiable and Analytic Manifolds, R, 9.3.7 that \(\phi_1\) and \(\phi_2\) coincide on a neighbourhood of \(e_{\mathrm{G}}\) .

COROLLARY 1. Let G and H be two Lie group germs. If L(G) and L(H) are isomorphic, G and H are locally isomorphic.

This follows from Theorem 1 and § 1, no. 10, Proposition 21.

COROLLARY 2. Let \(\mathbf{G}\) be a Lie group germ. If \(\mathrm{L}(\mathbf{G})\) is commutative, \(\mathbf{G}\) is locally isomorphic to the additive Lie group \(\mathrm{L}(\mathbf{G})\) .

The Lie algebra of the additive group \(\mathbf{L}(\mathbf{G})\) is isomorphic to \(\mathbf{L}(\mathbf{G})\) . Hence it suffices to apply Corollary 1.

COROLLARY 3. Let G be a Lie group. If L(G) is commutative, G contains a commutative open subgroup.

There exists an open Lie subgroup germ U of G which is commutative (Corollary 2). Let V be a neighbourhood of e such that \(V^{2} \subset U\) . Then xy = yx for all x, y in V. Hence the subgroup of G generated by V is commutative; it is obviously open.

\subsection{PASSAGE FROM LIE ALGEBRAS TO LIE GROUPS}

We shall denote the Hausdorff series by H(X, Y) (Chapter II, § 6, no. 4, Definition 1).

Lemma 2. Let \(\mathbf{L}\) be a complete normed Lie algebra over \(\mathbf{R}\) or \(\mathbf{C}\) . Let \(\mathbf{G}\) be the set of \(x \in \mathbf{L}\) such that \(\| x \| < \frac{1}{3} \log \frac{3}{2}\) . Let \(\theta\) be the mapping \(x \mapsto -x\) of \(\mathbf{G}\) into \(\mathbf{G}\) . Let \(\mathbf{H}\) be the restriction to \(\mathbf{G} \times \mathbf{G}\) of the Hausdorff function of \(\mathbf{L}\) (Chapter II, § 7, no. 2).

\begin{enumerate}
\def\labelenumi{(\roman{enumi})}
\item
  (G, 0, θ, H) is a Lie group germ.
\item
  Let \(\phi\) be the identity mapping of G into L. The differential of \(\phi\) at 0 is an isomorphism of the normable Lie algebra L(G) onto L.
\end{enumerate}
(i) follows from Chapter II, § 7, no. 2.

As \(\phi\) is a chart on G, the differential \(\psi\) of \(\phi\) at 0 is an isomorphism of normable spaces. On the other hand, the expansion as an integral series \(\mathrm{H} = \sum_{i,j\geqslant 0}\mathrm{H}_{ij}\) of the mapping H is such that \(\mathrm{H}_{11}(x,y) = \frac{1}{2} [x,y]\) . By § 3, Proposition 24, for all \(a,b\) in \(\mathrm{L}(\mathrm{G})\) ,

\[
\psi ([ a, b ]) = \mathrm{H} _ {1 1} (\psi (a), \psi (b)) - \mathrm{H} _ {1 1} (\psi (b), \psi (a)) = [ \psi (a), \psi (b) ]
\]

which proves (ii).

G is called the Lie group germ defined by L.

Suppose that K is ultrametric. Let p be the characteristic of the residue field of K. If \(p \neq 0\) , let \(\lambda = |p|^{1/(p-1)}\) ; if p = 0, let \(\lambda = 1\) .

Lemma 3. Let L be a complete normed Lie algebra over K. Let G be the set of \(x \in \mathbf{L}\) such that \(\| x \| < \lambda\) . Let H: G × G → G be the Hausdorff function of L (Chapter II, § 8, no. 3).

\begin{enumerate}
\def\labelenumi{(\roman{enumi})}
\item
  With the law of composition H, G is a Lie group in which 0 is identity element and -x the inverse of x for all \(x \in G\) .
\item
  Let \(\phi\) be the identity mapping of \(\mathbf{G}\) into \(\mathbf{L}\) . The differential of \(\phi\) at 0 is an isomorphism of the normable Lie algebra \(\mathrm{L}(\mathrm{G})\) onto \(\mathbf{L}\) .
\item
  For all \(\mu \in \mathbf{R}_+^*\) , let \(\mathrm{G}_{\mu}\) be the set of \(x \in \mathbf{L}\) such that \(\| x \| < \mu\) . Then the \(\mathrm{G}_{\mu}\) for \(\mu < \lambda\) , form a fundamental system of open and closed neighbourhoods of 0 and are subgroups of \(\mathbf{G}\) .
\end{enumerate}

Assertions (i) and (iii) follow from Chapter II, § 8, no. 3, Proposition 3 and (ii) can be proved as in Lemma 2.

G is called the Lie group defined by L.

THEOREM 2. Let \(\mathbf{L}\) be a complete normable Lie algebra. There exists a Lie group germ \(G\) such that \(\mathrm{L}(\mathbf{G})\) is isomorphic to \(\mathbf{L}\) . Two such Lie group germs are locally isomorphic.

The first assertion follows from Lemmas 2 and 3. The second assertion follows from Corollary 1 to Theorem 1 of no. 1.

COROLLARY 1. Let \(\mathbf{G}\) be a Lie group. There exists a neighbourhood of \(e\) which contains no finite subgroup distinct from \(\{e\}\) . If \(\mathbf{K} = \mathbf{R}\) or \(\mathbf{C}\) , there exists an open neighbourhood of \(e\) which contains no subgroup distinct from \(\{e\}\) .

We write \(\mathrm{L}(G) = \mathrm{L}\) . Choose a norm on \(\mathbf{L}\) defining the topology on \(\mathbf{L}\) and such that \(\| [x,y]\| \leqslant \| x\| \| y\|\) for all \(x,y\) in \(\mathbf{L}\) .

Suppose that \(\mathbf{K} = \mathbf{R}\) or \(\mathbf{C}\) . Let \(G'\) be the Lie group germ defined by L. There exist an open ball \(U'\) of centre 0 in \(G'\) and an isomorphism \(\phi\) of the Lie group germ \(U'\) onto an open neighbourhood \(U\) of \(e\) in G. Let \(V' = \frac{1}{2} U'\) , \(V = \phi(V')\) , H be a subgroup of G contained in V and \(h \in H\) . Let \(x = \phi^{-1}(h) \in V'\) . If \(x \neq 0\) , there exists an integer \(n > 0\) such that \(x, 2x, \ldots, nx\) are in \(V'\) , \((n + 1)x \in U'\) , \((n + 1)x \notin V'\) . Then \(h, h^2, \ldots, h^n\) are in V,

\[
h ^ {n + 1} \in \mathrm{U}, \quad h ^ {n + 1} \notin \mathrm{V},
\]

which is absurd. Hence \(\mathbf{H} = \{e\}\) .

Suppose that K is ultrametric. It suffices to prove the corollary when G is the Lie group associated with L. If \(g \in G\) , the powers of g evaluated in G are the elements of Zg evaluated in L. These are distinct if \(g \neq e\) . Hence G contains no finite subgroup distinct from \(\{e\}\) .

COROLLARY 2. Let \(k\) be a non-discrete closed subfield of \(\mathbf{K}\) , \(\mathbf{G}\) a Lie group over \(k\) and \(\mathbf{L} = \mathbf{L}(\mathbf{G})\) . Suppose that \(\mathbf{L}\) has a normable Lie \(\mathbf{K}\) -algebra structure \(\mathbf{L}'\) , compatible with the normable Lie \(k\) -algebra structure and invariant under the adjoint representation of \(\mathbf{G}\) . Then there exists on \(\mathbf{G}\) one and only one Lie \(\mathbf{K}\) -group structure compatible with the Lie \(k\) -group structure and for which the Lie algebra is \(\mathbf{L}'\) .

There exists a Lie group germ \(G_{1}\) over K such that \(\mathrm{L}(G_{1}) = \mathrm{L}'\) (Theorem 2). By Corollary 1 to Theorem 1 of no. 1, G and \(G_{1}\) , considered as Lie k-group germs, are locally isomorphic. Hence there exist an open neighbourhood \(G'\) of e in G and a Lie K-group germ structure on \(G'\) , with Lie algebra L, which is compatible with the Lie group germ structure over k. Let V be a symmetric open neighbourhood of e in G such that \(V^{2} \subset G'\) . Let \(g \in G\) . Then \(\phi = Int g\) is a k-isomorphism of a sufficiently small open Lie subgroup germ of \(G'\) onto an open Lie subgroup germ of \(G'\) ; and \(T_{e}(\phi)\) is K-linear and hence \(T_{x}(\phi)\) is K-linear for x sufficiently close to e; therefore the restriction of Int g to a sufficiently small open neighbourhood of e in V is K-analytic (Differentiable and Analytic Manifolds, R, 5.14.6). By §1, no. 9, Proposition 18, there exists on G an analytic K-manifold structure with which G is a Lie

K-group and V is an open sub-K-manifold of G. By translation, it is seen that the underlying k-manifold structure of G is the given structure. The Lie algebra of the Lie K-group G is the same as that of the open Lie K-subgroup germ V and hence is L'. Finally, the uniqueness stated in the corollary follows from § 3, no. 8, Proposition 32.

THEOREM 3. Let G be a Lie group germ and \(\mathfrak{h}\) a Lie subalgebra of \(\mathrm{L}(\mathrm{G})\) admitting a topological supplement. There exists a Lie subgroup germ H of G such that \(\mathrm{L}(\mathrm{H}) = \mathfrak{h}\) . If \(\mathrm{H}_1\) and \(\mathrm{H}_2\) are Lie subgroup germs of G such that

\[
\mathrm{L} (\mathrm{H} _ {1}) = \mathrm{L} (\mathrm{H} _ {2}) = \mathfrak {h},
\]

then \(\mathrm{H}_1\cap \mathrm{H}_2\) is open in \(\mathrm{H}_1\) and \(\mathrm{H}_2\)

There exists a Lie group germ H' with Lie algebra isomorphic to \(\mathfrak{h}\) (Theorem 2). Shrinking H' if necessary, it can be assumed that there exists a morphism \(\phi\) of H' into G such that \(\mathrm{L}(\phi)\) is an isomorphism of \(\mathrm{L}(\mathrm{H}')\) onto \(\mathfrak{h}\) (no. 1, Theorem 1). As \(\mathfrak{h}\) admits a topological supplement, \(\phi\) is an immersion at e. Hence, shrinking H' further, it can be assumed that \(\phi\) is an isomorphism of the manifold H' onto a submanifold of G. This proves the existence of H. The second assertion follows from the following proposition:

PROPOSITION 1. Let G be a Lie group germ and H and H' two Lie subgroup germs. In order that L(H) ⊃ L(H'), it is necessary and sufficient that H ∩ H' be open in H'. If H ∩ H' is open in H', then L(H') = L(H ∩ H') ⊂ L(H). Suppose that L(H) ⊃ L(H'). Let i, i' be the canonical injections of H, H' into G. By shrinking H' if necessary, it can be assumed that there exists a morphism ψ of H' into H such that L(ψ) is the canonical injection of L(H') into L(H) (no. 1, Theorem 1). Then L(i ∘ ψ) = L(i') and hence there exists a neighbourhood V of eH, in H' such that i ∘ ψ and i' coincide on V (Theorem 1). Therefore V ⊂ H, hence V ⊂ H ∩ H' and H ∩ H' is open in H' (§ 1, no. 10).

PROPOSITION 2. Let G be a Lie group over K, k a non-discrete closed subfield of K and H a Lie subgroup of the Lie k-group G. Suppose that L(H) is a vector sub-K-space of L(G) which admits a topological supplement. Then H is a Lie subgroup of the Lie K-group G.

There exists a Lie subgroup germ \(\mathrm{H}'\) of the Lie K-group G such that \(\mathrm{L}(\mathrm{H}') = \mathrm{L}(\mathrm{H})\) (Theorem 3). Consider G, H, H' as Lie \(k\) -group germs; Theorem 3 then proves that \(\mathrm{H} \cap \mathrm{H}'\) is open in H and H'. Hence there exists an open neighbourhood U of \(e\) in G such that \(\mathrm{U} \cap \mathrm{H}\) is a submanifold of G over K. Therefore, H is a Lie subgroup of the Lie K-group G ( \(\S 1\) , no. 3, Proposition 6).

\subsection{EXPONENTIAL MAPPINGS}

THEOREM 4. Let G be a Lie group germ, L its Lie algebra, V an open neighbourhood of 0 in L, \(\phi\) an analytic mapping of V into G such that \(\phi(0) = 0\) and \(\mathrm{T}_0(\phi) = \mathrm{Id}_L\) . The following conditions are equivalent:

\begin{enumerate}
\def\labelenumi{(\roman{enumi})}
\item
  For all \(b \in \mathbf{L}\) , \(\phi((\lambda + \lambda')b) = \phi(\lambda b)\phi(\lambda'b)\) for \(|\lambda|\) and \(|\lambda'|\) sufficiently small.
\item
  For all \(b \in \mathbf{L}\) and every integer \(n > 0\) , \(\phi_{*}(b^{n})\) is homogeneous of degree \(n\) in

\(\mathrm{U}(\mathrm{G})\) (\(\mathrm{T}_{0}^{(\infty)}(\mathrm{L})\) is identified with \(\mathrm{TS}(\mathrm{L})\) and \(b^{n}\) is evaluated in \(\mathrm{TS}(\mathrm{L})\)).

\item
  The mapping \(\phi_{*}\) of TS(L) into U(G) is compatible with the graduations of TS(L) and U(G).
\item
  The mapping \(\phi_*\) of TS(L) into U(G) is the canonical mapping of TS(L) into the enveloping algebra of L.
\item
  There exist a norm on \(\mathbf{L}\) defining the topology of \(\mathbf{L}\) and such that
\end{enumerate}

\[
\| [ x, y ] \| \leqslant \| x \| \| y \|
\]

for all \(x, y\) in \(\mathbf{L}\) and an open subgroup germ \(\mathbf{W} \subset \mathbf{V}\) of the Lie group germ defined by \(\mathbf{L}\) (no. 2) such that \(\phi | \mathbf{W}\) is an isomorphism of \(\mathbf{W}\) onto an open Lie subgroup germ of \(\mathbf{G}\) .

(v) \(\Rightarrow\) (i): obvious, for \((\lambda b).(\lambda^{\prime}b) = (\lambda +\lambda^{\prime})b\) in W for \(|\lambda |\) and \(|\lambda '|\) sufficiently small.
(i) \(\Rightarrow\) (ii): suppose that condition (i) is satisfied. Let \(b \in L\) . Let \(\psi\) be the restriction of \(\phi\) to \(V \cap Kb\) . By hypothesis, there exists a symmetric neighbourhood \(T\) of 0 in the additive Lie group \(Kb\) such that \(\psi|T\) is a morphism of the Lie group germ \(T\) into \(G\) . Hence

\[
\phi_ {*} (b ^ {n}) = (\psi | \mathrm{T}) _ {*} (b ^ {n}) = ((\psi | \mathrm{T}) _ {*} (b)) ^ {n} = (\phi_ {*} (b)) ^ {n},
\]

so that \(\phi_{*}(b^{n})\) is homogeneous of degree \(n\) in U(G).

(ii) \(\Rightarrow\) (iii): this follows from the fact that \(\mathsf{TS}^n (\mathbf{L})\) is the vector subspace of \(\mathsf{TS}(\mathbf{L})\) generated by the \(n\) -th powers of the elements of \(\mathbf{L}\) (Algebra, Chapter IV, § 5, Proposition 5).
(iii) \(\Rightarrow\) (iv): the canonical mapping of TS(L) into the enveloping algebra of L is the unique morphism of graded cogebras mapping 1 to 1 and extending \(\mathrm{Id}_{\mathrm{L}}\) (Chapter II, § 1, no. 5, Remark 3). Now \(\phi_{*}\) is a cogebra morphism and \(\phi_{*}|_{\mathrm{L}} = \mathrm{Id}_{\mathrm{L}}\) by hypothesis. If condition (iii) holds, it is seen that condition (iv) also holds.
(iv) \(\Rightarrow\) (v): suppose that condition (iv) is satisfied. Choose a norm on L defining the topology of L and such that \(\| [x,y]\| \leqslant \| x\| \| y\|\) for all \(x,y\) in L. Let H be the Lie group germ defined by the normed Lie algebra L. By Theorem 1, there exist an open subgroup germ \(S\subset V\) of H and an isomorphism \(\phi^{\prime}\) of S onto an open subgroup germ of G. As we know already that (v) \(\Rightarrow\) (iv), the mapping \(\phi_{*}\) of TS(L) into U(G) is the canonical mapping of TS(L) into the enveloping algebra of L. Thus \(\phi_{*}(t) = \phi_{*}^{\prime}(t)\) for all \(t\in T_0^{(\infty)}(L)\) . As \(\phi\) and \(\phi^{\prime}\) are analytic, \(\phi\) and \(\phi^{\prime}\) coincide on a neighbourhood of 0.

DEFINITION 1. Let G be a Lie group germ and L its Lie algebra. An exponential mapping of G is any analytic mapping \(\phi\) defined on an open neighbourhood of 0 in L, with values in G and satisfying the conditions of Theorem 4.

Theorem 4 implies immediately that, for every Lie group germ G, there exists an exponential mapping of G and that two exponential mappings of G coincide on a neighbourhood of 0.

Examples. (1) Let G be the additive group of a complete normable space E. The canonical isomorphism of L(G) onto E satisfies condition (i) of Theorem 4 and is therefore an exponential mapping of G.

\begin{enumerate}
\def\labelenumi{(\arabic{enumi})}
\setcounter{enumi}{1}
\tightlist
\item
  Let A be a complete normed unital associative algebra. Let \(\mathbf{A}^*\) be the Lie group consisting of the invertible elements of A. We identify \(\mathrm{L}(\mathrm{A}^*)\) with A (§3, no. 9, Corollary to Proposition 33). If \(\mathbf{K} = \mathbf{R}\) or \(\mathbf{C}\) , we know that the mapping exp of A into \(\mathbf{A}^*\) defined in Chapter II, §7, no. 3 satisfies condition (i) of Theorem 4 and hence is an exponential mapping. Now let K be ultrametric. Let \(p\) be the characteristic of the residue field of K. If \(p \neq 0\) , let \(\lambda = |p|^{1/(p-1)}\) ; if \(p = 0\) , let \(\lambda = 1\) . Let U be the set of \(x \in \mathbf{A}\) such that \(\|x\| < \lambda\) . We know (Chapter II, §8, no. 4) that the mapping exp of U into \(\mathbf{A}^*\) satisfies condition (i) of Theorem 4 and hence is an exponential mapping. Note that U is an additive subgroup of A.
\end{enumerate}

This example explains the terminology adopted in Definition 1.

Let \(G\) be a Lie group germ and \(\phi\) an exponential mapping of \(G\) . Then \(\phi\) is étale at 0 and hence there exists an open neighbourhood \(U\) of 0 in \(L(G)\) such that \(\phi(U)\) is open in \(G\) and \(\phi|U\) is an isomorphism of the analytic manifold \(U\) onto the analytic manifold \(\phi(U)\) .

A canonical chart (of the first species) on G is a chart \(\psi\) on the analytic manifold G whose inverse mapping is an exponential mapping. If further G is finite-dimensional and a basis of L(G) is chosen, the coordinate system defined by \(\psi\) and this basis in the domain of \(\psi\) is called a canonical coordinate system (of the first species).

PROPOSITION 3. Let \(\mathbf{G}\) be a Lie group germ, \(\mathbf{L}\) its Lie algebra and \(\phi\) an exponential mapping of \(\mathbf{G}\) . Let \(\mathbf{L}_1, \ldots, \mathbf{L}_n\) be vector subspaces of \(\mathbf{L}\) such that \(\mathbf{L}\) is the topological direct sum of \(\mathbf{L}_1, \ldots, \mathbf{L}_n\) . The mapping

\[
(b _ {1}, b _ {2}, \dots , b _ {n}) \mapsto \theta (b _ {1}, b _ {2}, \dots , b _ {n}) = \phi (b _ {1}) \phi (b _ {2}) \dots \phi (b _ {n}),
\]

defined on an open subset of \(L_{1} \times L_{2} \times \cdots \times L_{n}\) , is analytic. The tangent mapping at \((0, 0, \ldots, 0)\) to \(\theta\) is the canonical mapping of \(L_{1} \times \cdots \times L_{n}\) into L.

Let \(k_{i}\) be the canonical injection of \(L_{i}\) into \(L_{1} \times L_{2} \times \cdots \times L_{n}\) . Then, for all \(b \in L_{i}\) , \((\mathrm{T}_{(0,\ldots,0)}\theta)(\mathrm{T}_{0}\phi_{i})(b) = (\mathrm{T}_{0}\phi)(b) = b\) and hence \((\mathrm{T}_{(0,\ldots,0)}\theta)|L_{i}\) is the canonical injection of \(L_{i}\) into L.

In particular, \(\theta\) is étale at \((0, 0, \ldots, 0)\) . Its restriction to a sufficiently small open neighbourhood U of \((0,0,\ldots,0)\) has an open image in G and is an isomorphism of the manifold U onto the manifold \(\theta(\mathrm{U})\) . The inverse mapping \(\eta\) of \(\theta(\mathrm{U})\) onto U is called a canonical chart of the second species of G, associated with the given decomposition of L as a direct sum. If further G is finite-dimensional and each \(L_{i}\) is generated by a non-zero vector \(e_{i}\) , the coordinate system on \(\theta(\mathrm{U})\) defined by \(\eta\) and the \(e_{i}\) is called a canonical coordinate system of the second species.

PROPOSITION 4. Let G be a Lie group germ and \(\phi\) an injective exponential mapping of G. For all \(x, y\) in L(G),


\[
x + y = \lim _ {\lambda \in \mathbb {K} ^ {*}, \lambda \rightarrow 0} \lambda^ {- 1} \phi^ {- 1} (\phi (\lambda x) \phi (\lambda y))\tag{1}
\]

\[
[ x, y ] = \lim _ {\lambda \in \mathbb {K} ^ {*}, \lambda \rightarrow 0} \lambda^ {- 2} \phi^ {- 1} (\phi (\lambda x) \phi (\lambda y) \phi (- \lambda x) \phi (- \lambda y))\tag{2}
\]

(note that \(\phi^{-1}(\phi(\lambda x)\phi(\lambda y))\) and \(\phi^{-1}(\phi(\lambda x)\phi(\lambda y)\phi(-\lambda x)\phi(-\lambda y))\) are defined for \(|\lambda|\) sufficiently small).

Let \(\mathbf{L} = \mathbf{L}(\mathbf{G})\) be given a norm defining the topology of \(\mathbf{L}\) and such that \(\| [x,y] \| \leqslant \| x \| \| y \|\) for all \(x, y\) in \(\mathbf{L}\) . Using Theorems 2 and 4, it can be assumed that \(\mathbf{G}\) is the Lie group germ defined by \(\mathbf{L}\) and that \(\phi = \mathrm{Id}_{\mathbf{G}}\) . Let \((x,y) \mapsto x.y\) denote the product in the group \(\mathbf{G}\) . The formulae to be proved can then be written


\[
x + y = \lim _ {\lambda \in \mathbf {K} ^ {*}, \lambda \rightarrow 0} \lambda^ {- 1} ((\lambda x). (\lambda y))\tag{3}
\]

\[
[ x, y ] = \lim _ {\lambda \in \mathbb {K} ^ {*}, \lambda \rightarrow 0} \lambda^ {- 2} ((\lambda x). (\lambda y). (- \lambda x). (- \lambda y)).\tag{4}
\]

There exists an open neighbourhood V of 0 in K such that the function

\[
\lambda \mapsto f (\lambda) = (\lambda x). (\lambda y)
\]

is defined and analytic on V. By Chapter II, § 6, no. 4, Remark 2, the expansion of \(f\) as an integral series about the origin is

\[
\lambda (x + y) + \frac {1}{2} \lambda^ {2} [ x, y ] + \dots
\]

and this proves (3). On the other hand, for \(u, v\) in G and \(\| u \|\) , \(\| v \|\) sufficiently small, \(u.v\) is an analytic function of \((u, v)\) and the terms of degrees 1 and 2 in the expansion of this function as an integral series about the origin are \(u + v + \frac{1}{4}[u, v]\) . By Differentiable and Analytic Manifolds, R, 3.2.7 and 4.2.3, the terms of degrees 1 and 2 in the expansion of the function \(f(\lambda)f(-\lambda)\) as an integral series about the origin are the terms of degrees 1 and 2 in

\[
f (\lambda) + f (- \lambda) + \frac {1}{2} [ f (\lambda), f (- \lambda) ]
\]

or also in

\[
\begin{array}{r l} \lambda (x + y) + \frac {1}{2} \lambda^ {2} [ x, y ] - \lambda (x + y) + \frac {1}{2} \lambda^ {2} [ x, y ] & \\ + \frac {1}{2} [ \lambda (x + y), - \lambda (x + y) ] = \lambda^ {2} [ x, y ] \end{array}
\]

and this proves (4).

PROPOSITION 5. Let G be a Lie group, k a non-discrete closed subfield of K, G' the group G considered as a Lie group over k and \(\phi\) (resp. \(\phi'\) ) an exponential mapping of G (resp. G'). Then \(\phi\) and \(\phi'\) coincide on a neighbourhood of 0.

\(\phi\) satisfies hypothesis (i) of Theorem 4 relative to \(G'\) and is therefore an exponential mapping of \(G'\) .

PROPOSITION 6. Let G be a Lie group germ, L its Lie algebra and \(\phi: V \to G\) an exponential mapping of G. For all \(x \in V\) , let \(\mathrm{T}_x(L)\) be identified with L, so that the right differential \(\varpi(x)\) of \(\phi\) at \(x\) is a linear mapping of L into L. For \(x\) sufficiently close to 0,

\[
\varpi (x) = \sum_ {n \geqslant 0} \frac {1}{(n + 1) !} (\mathrm{ad} x) ^ {n}.
\]

Let \(\mathbf{L}\) be given a norm compatible with its topology and such that \(\| [x,y]\| \leqslant \| x\| \| y\|\) for all \(x,y\in L\) . It suffices to consider the case where \(G\) is the Lie group germ defined by \(L\) and \(\phi = \mathrm{Id}_{\mathrm{G}}\) . By definition, \(\varpi (x)\) is then the tangent mapping at \(x\) to the mapping \(y\mapsto y.x^{-1}\) of \(G\) into \(G\) . If \(H(X,Y)\) denotes the Hausdorff series, \(\varpi (x)\) is therefore, for \(\| x\|\) sufficiently small, the tangent mapping at 0 to the mapping \(y\mapsto H(x + y, - x)\) of \(G\) into \(G\) . In \(H(X + Y, - X)\) , the sum of terms of the first degree in \(Y\) is

\[
\sum_ {m \geqslant 0} \frac {1}{(m + 1) !} (\mathrm{adX}) ^ {m} \mathrm{Y}
\]

(Chapter II, § 6, no. 5, Proposition 5). The proposition then follows from Differentiable and Analytic Manifolds, R, 3.2.4 and 4.2.3.

Let G be a Lie group germ and \(t \in K\) . A t-th power mapping of G is any mapping, defined and analytic on an open neighbourhood of e, with values in G and coinciding on a neighbourhood of e with a mapping

\[
g \mapsto \phi (t \phi^ {- 1} (g))
\]

where \(\phi\) is an injective exponential mapping of G.

PROPOSITION 7. (i) If \(t \in \mathbf{Z}\) , a \(t\) -th power mapping coincides on a neighbourhood of \(e\) with the mapping \(g \mapsto g^t\) .

\begin{enumerate}
\def\labelenumi{(\roman{enumi})}
\setcounter{enumi}{1}
\item
  The tangent mapping at \(e\) to a \(t\) -th power mapping is the homothety of ratio \(t\) .
\item
  If \(h\) is a \(t\) -th power mapping and \(h'\) a \(t'\) -th power mapping of \(G\) , \(h \circ h'\) is a \((tt')\) -th power mapping and \(g \mapsto h(g)h'(g)\) is a \((t + t')\) -th power mapping.
\item
  If \(h\) is a \(t\) -th power mapping and \(u \in \mathbf{U}^n(\mathbf{G})\) , then \(h_*(u) = t^n u\) .
\end{enumerate}

It suffices to prove the proposition when G is the Lie group germ defined by a complete normed Lie algebra and when the t-th power mappings considered are constructed using the exponential mapping \(\phi = \mathrm{Id}_{\mathrm{G}}\) . But in that case everything is obvious.

\subsection{FUNCTORIALITY OF EXPONENTIAL MAPPINGS}

PROPOSITION 8. Let G and H be Lie group germs, h a morphism of G into H and \(\phi_{\mathrm{G}}\) and \(\phi_{\mathrm{H}}\) exponential mappings of G and H. There exists a neighbourhood V of 0 in L(G) such that \(h \circ \phi_{\mathrm{G}}\) and \(\phi_{\mathrm{H}} \circ \mathrm{L}(h)\) coincide on V.

Let \(\mathrm{L}(\mathrm{G})\) and \(\mathrm{L}(\mathrm{H})\) be given norms defining their topologies and such that \(\| [x,y]\| \leqslant \| x\| \| y\|\) for all \(x\) and \(y\) . It can be assumed that \(\mathrm{G}\) (resp. H) is the Lie group germ defined by \(\mathrm{L}(\mathrm{G})\) (resp. \(\mathrm{L}(\mathrm{H})\) ), so that \(\phi_{\mathrm{G}}\) (resp. \(\phi_{\mathrm{H}}\) ) coincides with \(\mathrm{Id}_{\mathrm{G}}\) (resp. \(\mathrm{Id}_{\mathrm{H}}\) ) in a neighbourhood of 0. On the other hand, there exists a symmetric open neighbourhood W of 0 in \(\mathrm{L}(\mathrm{G})\) such that \(\mathrm{L}(h)\) is a morphism of the Lie group germ W into H. By Theorem 1, \(\mathrm{L}(h)\) coincides with \(h\) on a neighbourhood of 0, whence the proposition.

Loosely speaking, if G and H are identified in a neighbourhood of the identity element with L(G) and L(H) by means of exponential mappings, every morphism of G into H is linear in a neighbourhood of 0.

COROLLARY 1. Let G be a Lie group germ, G' a Lie subgroup germ of G and \(\phi\) an exponential mapping of G.

\begin{enumerate}
\def\labelenumi{(\roman{enumi})}
\item
  There exists an open neighbourhood \(\mathbf{V}\) of 0 in \(\mathrm{L}(\mathbf{G}')\) such that \(\phi | \mathbf{V}\) is an isomorphism of the manifold \(\mathbf{V}\) onto an open neighbourhood of \(e\) in \(\mathbf{G}'\) .
\item
  Let \(x \in \mathrm{L}(\mathrm{G})\) . The following conditions are equivalent: (a) \(x \in \mathrm{L}(\mathrm{G}')\) ; (b) \(\phi(\lambda x) \in \mathrm{G}'\) for \(|\lambda|\) sufficiently small.
\item
  is obtained by applying Proposition 8 to the canonical injection of \(\mathbf{G}'\) into \(\mathbf{G}\) and (ii) follows from (i).
\end{enumerate}

COROLLARY 2. Let G be a Lie group, \(\rho\) an analytic linear representation of G and \(\phi\) an exponential mapping of G. There exists a neighbourhood V of 0 in L(G) such that

for all \(x \in V\) .

\[
\rho (\phi (x)) = \exp (\mathrm{L} (\rho) x)
\]

This follows from Proposition 8 and Example 2 of no. 3.

COROLLARY 3. Let G be a Lie group and \(\phi\) an exponential mapping of G.

\begin{enumerate}
\def\labelenumi{(\roman{enumi})}
\tightlist
\item
  There exists a neighbourhood V of 0 in L(G) such that
\end{enumerate}

\[
\operatorname{Ad} (\phi (x)) = \exp \operatorname{ad} x.
\]

for all \(x \in V\) .

\begin{enumerate}
\def\labelenumi{(\roman{enumi})}
\setcounter{enumi}{1}
\tightlist
\item
  If \(g \in \mathbf{G}\) , there exists a neighbourhood \(\mathbf{W}\) of 0 in \(\mathrm{L}(\mathbf{G})\) such that
\end{enumerate}

\[
g \phi (x) g ^ {- 1} = (\mathrm{Ad} g. x)
\]

for all \(x \in W\) .

\begin{enumerate}
\def\labelenumi{(\roman{enumi})}
\tightlist
\item
  follows from Corollary 2 and § 3, no. 12, Proposition 44.
\end{enumerate}

\[
g.
\]

\subsection{STRUCTURE INDUCED ON A SUBGROUP}

Lemma 4. Let G be a finite-dimensional Lie group, \(\Omega\) a symmetric open neighbourhood of \(e\) in G and H a subset of \(\Omega\) containing \(e\) such that the conditions \(x \in \mathrm{H}\) , \(y \in \mathrm{H}\) , \(xy^{-1} \in \Omega\) imply \(xy^{-1} \in \mathrm{H}\) . Let \(r \in \mathrm{N}_{\mathbb{K}}\) . For all \(x \in \mathrm{H}\) , let \(\mathfrak{h}_x\) be the set of \(a \in \mathrm{T}_x(\mathrm{G})\) with the following property: there exist an open neighbourhood I of 0 in K and a mapping \(f\) of class \(\mathrm{C}^r\) of I into G such that \(f(0) = x, f(I) \subset H\) , \((\mathrm{T}_0f)(1) = a\) .

\begin{enumerate}
\def\labelenumi{(\roman{enumi})}
\item
  Let \(\mathfrak{h}_e = \mathfrak{h}\) . Then \(\mathfrak{h}\) is a Lie subalgebra of \(\mathrm{L}(\mathrm{G})\) which is invariant under \(\mathrm{Ad}_{\mathrm{L}(\mathrm{G})}(\mathrm{H})\) .
\item
  \(\mathfrak{h}_x = x\mathfrak{h} = \mathfrak{h}x\) for all \(x\in \mathrm{H}\) (where \(x\mathfrak{h}\) and \(\mathfrak{h}x\) are evaluated in \(\mathrm{T(G)}\) ).
\item
  Let \(\mathbf{V}\) be a manifold of class \(\mathbf{C}^r\) , \(v_0\) a point of \(\mathbf{V}\) and \(f\) a mapping of class \(\mathbf{C}^r\) of \(\mathbf{V}\) into \(\mathbf{G}\) such that \(f(v_0) = e\) and \(f(\mathbf{V}) \subset \mathbf{H}\) . For every Lie subgroup germ \(\mathbf{H}'\) of \(\mathbf{G}\) with Lie algebra \(\mathfrak{h}, f(v) \in \mathbf{H}'\) for \(v\) sufficiently close to \(v_0\) .
\item
  For every Lie subgroup germ \(\mathbf{H}'\) of \(\mathbf{G}\) with Lie algebra \(\mathfrak{h}\) , \(\mathbf{H}' \cap \mathbf{H}\) is a neighbourhood of \(e\) in \(H'\) .
\item
  For all \(x \in \mathbf{H}\) and all \(a \in \mathfrak{b}_x\) , there exist an open neighbourhood I of 0 in K and a mapping \(f\) of class \(\mathbf{C}^{\omega}\) of I into G such that \(f(0) = x, f(I) \subset H, (T_0f)(1) = a\) . Clearly \(\mathrm{K}\mathfrak{b} = \mathfrak{b}\) and \(x\mathfrak{b}_y z = \mathfrak{b}_{xyz}\) for \(x, y, xy, xyz\) in H. This implies (ii) and the fact that \(\mathfrak{b}\) is invariant under \(\mathrm{Ad}_{\mathrm{L(G)}}(\mathrm{H})\) .
\end{enumerate}

Let \(a_1, a_2\) be in \(\mathfrak{h}\) . Let I be an open neighbourhood of 0 in K and \(f_1, f_2\) mappings of class \(\mathrm{C}^r\) of I into G such that \(f_j(0) = e, f_j(I) \subset H\) , \((\mathrm{T}_0 f_j)(1) = a_j (j = 1, 2)\) . We define \(f: I \to G\) by \(f(\lambda) = f_1(\lambda)f_2(\lambda)\) . Then \(f\) is of class \(\mathrm{C}^r\) and \(f(0) = e\) . By shrinking I if necessary, we have \(f(I) \subset H\) . On the other hand, the mapping of \(\mathrm{T}_e(G) \times \mathrm{T}_e(G)\) into \(\mathrm{T}_e(G)\) tangent to the mapping \((g, g') \to gg'\) is addition; hence \((\mathrm{T}_0 f) = a_1 + a_2\) . Hence \(a_1 + a_2 \in \mathfrak{h}\) and \(\mathfrak{h}\) is a vector subspace of L(G). Since \(x\mathfrak{h}x^{-1} = \mathfrak{h}\) for all \(x \in H\) , \((\mathrm{Ad}f_1(\lambda)).a_2 \in \mathfrak{h}\) for all \(\lambda \in I\) . The tangent mapping at 0 to the mapping \(\lambda \mapsto \mathrm{Ad}f_1(\lambda)\) is, by Proposition 44 of § 3, no. 12, the mapping \(\lambda \mapsto \mathrm{ad}(\lambda a_1)\) ; hence

\[
\left[ a _ {1}, a _ {2} \right] = (\operatorname{ad} a _ {1}). a _ {2} \in \mathfrak {h}
\]

since \(\mathfrak{h}\) is closed in \(\mathrm{L}(\mathrm{G})\) . Hence we have proved (i). In the rest of the proof we fix a Lie subgroup germ \(\mathrm{H}'\) of \(\mathrm{G}\) with Lie algebra \(\mathfrak{h}\) .

Let \(V, v_0, f\) be as in (iii). Let \(Y\) be the left foliation of \(G\) associated with \(\mathfrak{h}\) (no. 1). For all \(y \in H'\) , \(T_y(H') = y\mathfrak{h}\) . On the other hand, for all \(v \in V\) , the image of \(T_v(V)\) under \(T_v(f)\) is contained in \(\mathfrak{h}_{f(v)} = f(v)\mathfrak{h}\) (by definition of \(\mathfrak{h}_{f(v)}\) ). By Differentiable and Analytic Manifolds, R, 9.3.2, \(f\) is a morphism of \(V\) into \(Y\) . As \(H'\) is a leaf of \(Y\) (Differentiable and Analytic Manifolds, R, 9.2.8), \(f(v) \in H'\) for \(v\) sufficiently close to \(v_0\) .

Let \((a_{1},\ldots ,a_{s})\) be a basis of \(\mathfrak{h}\) . There exist an open neighbourhood I of 0 in K and mappings \(f_{1},\ldots ,f_{s}\) of class \(\mathbf{C}^r\) of I into G such that \(f_{j}(0) = e\) , \(f_{j}(\mathrm{I})\subset \mathrm{H}\) , \((\mathrm{Tf}_j)1 = a_j\) for all \(j\) . By (iii), \(f_{j}(\lambda)\in \mathrm{H}'\) for \(|\lambda |\) sufficiently small. Hence the \(f_{1}(\lambda_{1})f_{2}(\lambda_{2})\ldots f_{s}(\lambda_{s})\) constitute, for \(|\lambda_1|,\left|\lambda_2\right|,\ldots ,|\lambda_s|\) sufficiently small, a neighbourhood of \(e\) in \(\mathbf{H}'\) ; and this neighbourhood is contained in \(\mathbf{H}\) . Hence (iv).

If \(a \in \mathfrak{h}\) , there exist an open neighbourhood \(\mathbf{I}\) of 0 in \(\mathbf{K}\) and a mapping \(f\) of class \(\mathbf{C}^{\omega}\) of \(\mathbf{I}\) into \(\mathbf{G}\) such that \(f(0) = e, f(\mathbf{I}) \subset \mathbf{H}'\) , \((\mathrm{T}_0f)1 = a\) . This, with (iv), implies (v).

DEFINITION 2. \(\mathfrak{h}\) is called the subalgebra tangent to H at \(e\) .

PROPOSITION 9. Let G be a finite-dimensional Lie group and H a subgroup of G.

\begin{enumerate}
\def\labelenumi{(\roman{enumi})}
\item
  There exists on H one and only one analytic manifold structure with the following property: for all r between 1 and \(\omega\) , for every manifold V of class \(C^r\) and for every mapping f of V into H, \(f\) is of class \(C^r\) as a mapping of V into H if and only if f is of class \(C^r\) as a mapping of V into G.
\item
  With this structure, \(\mathbf{H}\) is a Lie group, the canonical injection \(i\) of \(\mathbf{H}\) into \(\mathbf{G}\) is an immersion and \(\mathbf{L}(i)(\mathbf{L}(\mathbf{H}))\) is the Lie subalgebra tangent at \(e\) to \(\mathbf{H}\) .
\end{enumerate}

In (i) the uniqueness is obvious. We prove the existence. Let \(\mathfrak{b}\) be the Lie algebra tangent at \(e\) to H. Let H' be a Lie subgroup germ of G with Lie algebra \(\mathfrak{b}\) . By replacing H' by an open subgroup germ of H', it can be assumed that H' ⊆ H (Lemma 4 (iv)). For all \(x \in \mathrm{H}\) , \(xH'x^{-1}\) is a Lie subgroup germ of G with Lie algebra \(x \mathfrak{h}x^{-1} = \mathfrak{b}\) . Hence H' ∩ \((xH'x^{-1})\) is open in H' (no. 2, Theorem 3) and the mapping y → \(xyx^{-1}\) is an isomorphism of H' ∩ \(x^{-1}H'x\) onto \(xH'x^{-1}\) ∩ H'. Using Proposition 18 of § 1, no. 9, there exist an open Lie subgroup germ W of H' and a Lie group structure on H with the following properties: W is open in H and the manifold structures on H and H' induce the same structure on W. It then follows that the canonical injection i of H into G is an immersion and that L(i)(L(H)) = L(H') = \(\mathfrak{b}\) . Moreover, let V and f be as in (i). If f: V → H is of class \(C^r\), \(i \circ f\): V → G is of class \(C^r\). Suppose that \(i \circ f\): V → G is of class \(C^r\); then we prove that f: V → H is of class \(C^r\). By translation, it suffices to consider the case where there exists v0 ∈ V such that f(v0) = e and to prove that f: V → H is of class \(C^r\) in an open neighbourhood of v0. Now, by Lemma 4 (iii), f(v) ∈ H' for v sufficiently close to v0, whence our assertion. Thus we have proved (i) and (ii) has been obtained on the way.

DEFINITION 3. The Lie group structure on H defined in Proposition 9 is called the structure induced on H by the Lie group structure on G.

If H is a Lie subgroup of G, its Lie group structure is induced by that on G (Differentiable and Analytic Manifolds, R, 5.8.5).

If \(\mathbf{G} = \mathbf{R}\) and \(\mathbf{H} = \mathbf{Q}\) , then \(\mathfrak{h} = \{0\}\) and hence the structure induced on \(\mathbf{H}\) is the discrete Lie group structure. Similarly if \(\mathbf{G} = \mathbf{C}\) (considered as a complex Lie group) and \(\mathbf{H} = \mathbf{R}\) .
